\section{Concrete axioms for $A_\infty$ chiral algebras}
\label{sec:axioms-Ainfty-chiral}

An $\Ainf$ chiral algebra over $C_\ast(W(\SCchtop))$ is operadically
defined but lacks a finite list of structure maps: the Stasheff
relations sit inside the Boardman--Vogt resolution rather than as
identities among named operations. A coordinate presentation by
multilinear operations $\{m_k\}$ with relations as polynomial
identities in $\lambda$-brackets, verifiable case by case
(Virasoro, affine Kac--Moody, $\beta\gamma$), is what a vertex
algebraist computes.

The axiomatic list is equivalent to the operadic one.
Sesquilinearity encodes holomorphic translation covariance; the
Stasheff identities are Stokes' theorem on the Fulton--MacPherson
compactification $\FM_k(\C)$; the planar ordering convention
records the $E_1$ colour on~$\R$.
Theorem~\ref{thm:rectification} and
Proposition~\ref{prop:operad-implies-axioms} give the coordinate
expression of the operad.

\subsection{Underlying data and parameters}
\label{subsec:underlying-data}

The minimal data presenting an $\Ainf$ chiral algebra concretely is:
a cochain complex $(A, Q)$, a holomorphic translation operator
$\partial$, a vacuum unit $\mathbf{1}$, and a collection $\{m_k\}$ of
multilinear operations valued in iterated Laurent series in spectral
parameters~$\lambda_i$. Sesquilinearity, $A_\infty$ relations, and
unitality follow in the subsequent subsections; the present
subsection fixes only the underlying data.

\begin{definition}[$A_\infty$ chiral algebra]
\label{def:ainfty_chiral}
Fix a ground field $k$ of characteristic $0$. Let $(A,Q)$ be a cochain complex (cohomological grading), with $Q$ of degree $+1$.
Choose a closed vacuum unit $\mathbf 1\in A^0$ and put
$\bar A:=A/k\cdot\mathbf 1$. The first object is the coaugmented
chiral tensor coalgebra and its coaugmentation coideal
\[
\mathrm B^{\mathrm{ch}}(A):=T^c(s^{-1}\bar A)
=\bigoplus_{n\ge0}(s^{-1}\bar A)^{\otimes n},
\qquad
\barBch(A):=\overline{\mathrm B}^{\mathrm{ch}}(A)
=\bigoplus_{n\ge1}(s^{-1}\bar A)^{\otimes n}.
\]
The $n=0$ summand is $k$, and all coderivations below preserve
$\barBch(A)$. We equip $A$ with:
\begin{itemize}
 \item a degree $0$ endomorphism $\partial\colon A\to A$ (holomorphic translation) with $[Q,\partial]=0$,
 \item the chosen unit satisfying $Q\mathbf 1=0$ and $\partial \mathbf 1=0$,
 \item multilinear operations (for $k\ge 1$)
 \begin{equation}
 \label{eq:mk-type}
 m_k \colon A^{\otimes k}\longrightarrow A((\lambda_1))\cdots((\lambda_{k-1})),
 \qquad \deg m_k = 2-k,
 \end{equation}
 where $A((\lambda))$ denotes the space of formal Laurent series in $\lambda$ with coefficients in $A$, and $A((\lambda_1))\cdots((\lambda_{k-1}))$ is the iterated Laurent series ring.\footnote{Work in the completed topology so that infinite sums are allowed.}
\end{itemize}
We regard $\lambda_1,\dots,\lambda_{k-1}$ as spectral parameters associated to the \emph{differences} of the holomorphic coordinates of inputs $1,\dots,k-1$ relative to the $k$-th input (the ``rightmost'' slot).
The unary operation is $m_1=Q$. For a homogeneous input $a_i$, put
$\bar a_i:=|a_i|-1=|s^{-1}a_i|$. The operations above are the
cogenerator projections of a degree-$1$ coderivation $D_A$ of
$\mathrm B^{\mathrm{ch}}(A)$:
\[
\pi D_A\bigl[s^{-1}a_1|\cdots|s^{-1}a_k\bigr]
=s^{-1}m_k(a_1,\ldots,a_k),
\qquad
\pi\colon\barBch(A)\to s^{-1}\bar A .
\]
On a homogeneous bar word,
\begin{multline}
\label{eq:axioms-ainfty-coderivation}
D_A\bigl[s^{-1}a_1|\cdots|s^{-1}a_n\bigr] \\
=
\sum_{\substack{k\ge1\\1\le i\le n-k+1}}
(-1)^{\epsilon(i,k)}
\bigl[s^{-1}a_1|\cdots|
s^{-1}m_k(a_i,\ldots,a_{i+k-1})|\cdots|s^{-1}a_n\bigr],
\qquad
\epsilon(i,k)=\sum_{r<i}\bar a_r .
\end{multline}
The deconcatenation coproduct obeys
$\Delta D_A=(D_A\otimes\id+\id\otimes D_A)\Delta$ on
$\mathrm B^{\mathrm{ch}}(A)$. The single equation
$D_A^2=0$, expanded below with the same suspension convention, is
the $\Ainf$ chiral relation.
The axioms imposed on this data are stated in Definition~\ref{def:sesquilinearity}--\ref{def:symmetry-grading} below.
\end{definition}

\begin{remark}[Categorical origin of the $A_\infty$ chiral algebra]
\label{rem:categorical-origin}
The $A_\infty$ chiral algebra $A$ above is the endomorphism algebra $A_b = \RHom_{\cC(J)}(b,b)$ of a vacuum object $b$ in the open-sector factorization category $\cC(J)$ attached to a bulk chiral algebra $J$ (Definition~\ref{def:oc-factorization-category}). The multilinear operations $m_k$ arise from restricting the factorization functors of $\cC(J)$ to iterated self-$\Hom$'s of~$b$ (Theorem~\ref{thm:boundary-algebra-from-vacuum}). The complete invariant of the open sector is the Morita class of~$\cC(J)$, not the quasi-isomorphism class of~$A_b$: distinct compact generators $b, b'$ produce Morita-equivalent algebras $A_b \sim_M A_{b'}$ with equivalent module categories.
\end{remark}

\begin{remark}[Five primitive objects]
\label{rem:five-primitive-objects-axioms}
The boundary chart $A=A_b$ is one member of the primitive package.
The coaugmented chiral bar coalgebra
$\mathrm B^{\mathrm{ch}}(A)=k\oplus\barBch(A)$ resolves twisting
morphisms out of~$A$, with reduced coideal
$\barBch(A)=\bigoplus_{n\ge1}(s^{-1}\bar A)^{\otimes n}$ carrying the
bar codifferential. The linear Koszul dual $A^!$ is the cobar algebra built
from the dual cooperadic data; it differs from $\barB^{\mathrm{ch}}(A)$
and from the coaugmentation coideal $A^i$. The derived chiral centre
$\cZ^{\mathrm{der}}_{\mathrm{ch}}(A)$ is the Hochschild object
controlling bulk observables acting on the boundary. The operad
$\SCchtop$ is none of these: it is the two-coloured governing
operad whose open colour acts on~$A$ and whose closed colour acts
on $\cZ^{\mathrm{der}}_{\mathrm{ch}}(A)$.
\end{remark}

Definition~\ref{def:ainfty_chiral} is the zero-curvature case: $A$
sits over a point. Families over a nontrivial base $S$ (genus-$g$
moduli $\overline{\cM}_g$, level space, coupling parameter)
generically carry an obstruction to $m_1^2 = 0$, detected by a
degree-$2$ section of $\cO_S$; the natural enlargement absorbs the
obstruction as a curvature element $m_0$ of the $A_\infty$
structure.

\begin{definition}[Curved $A_\infty$ chiral algebra]
\label{def:curved-ainfty-chiral}
A \emph{curved $A_\infty$ chiral algebra} over a base $S$ (e.g., $S = \overline{\cM}_g$ for genus-$g$ modular theory) is the data of Definition~\ref{def:ainfty_chiral} together with:
\begin{itemize}
\item a \emph{curvature element} $m_0 \in \Gamma(S, \cO_S)$ of cohomological degree~$2$ (a section of the structure sheaf of the base, with no algebra inputs),
\end{itemize}
subject to the \emph{curved $A_\infty$ relations}: for each $n \geq 0$,
\[
\sum_{\substack{i+j = n+1 \\ i,j \geq 0}} \sum_{s=0}^{n-j} (-1)^{\epsilon(s,j)}\, m_i\bigl(\ldots, m_j(\ldots), \ldots\bigr) \;=\; 0,
\]
where the $j=0$ terms contribute $m_0$ insertions. The first three cases are:
\begin{itemize}
\item \textbf{At $n=0$:}\; $m_1(m_0) = 0$. The curvature element $m_0$ is closed under the differential~$m_1$.
\item \textbf{At $n=1$:}\; $m_1^2(x) = m_2(m_0, x) - m_2(x, m_0) = [m_0, x]_{m_2}$. The differential does not square to zero; rather, it squares to the commutator with the curvature (the graded commutator reduces to the ordinary commutator because $|m_0| = 2$ is even).
\item \textbf{At $n=2$:}\; $m_1(m_2(x,y)) + m_2(m_1(x),y) + (-1)^{|x|}\,m_2(x,m_1(y)) + m_3(m_0,x,y) \pm m_3(x,m_0,y) \pm m_3(x,y,m_0) = 0$. The binary product is $m_1$-linear up to homotopies controlled by $m_3$ and curvature insertions (setting $m_0=0$, this is the graded Leibniz rule: $m_1$ is a derivation of~$m_2$).
\end{itemize}
When $m_0 = \kappaChHodge(\cA) \cdot \omega_g$ for a scalar $\kappa$ and a base class $\omega_g$, the $n=1$ relation reads $d_{\barB}^2 = \kappa \cdot \omega_g$ (Volume~I, Theorem~D): the bar differential squares to the curvature obstruction, and this identity is the non-formality at genus $g \geq 1$. Setting $m_0 = 0$ recovers Definition~\ref{def:ainfty_chiral}.
\end{definition}

\begin{notation}[Spectral convention]
\label{not:spectral}
For $n\ge1$ and ordered inputs $a_1,\dots,a_n$, we always write $m_n(a_1,\dots,a_n)$ as a formal Laurent series in $(\lambda_1,\dots,\lambda_{n-1})$ following the convention above. If $I=[p+1,\dots,p+r]$ is a consecutive block, we denote
\[
\Lambda_I\;:=\;\lambda_{p+1}+\cdots+\lambda_{p+r-1},
\]
and for any index $j>p+r-1$ we define $\lambda^{(I)}_j:=\lambda_j$ (unchanged), whereas the \emph{effective parameter} for the collapsed block $I$ is $\Lambda_I$.
\end{notation}

\begin{remark}[Operadic formulation]
\label{rem:operadic-formulation}
\index{chiral endomorphism operad!operadic formulation}
The data of Definition~\ref{def:ainfty_chiral} admits an operadic
reformulation. The \emph{chiral endomorphism operad}
$\cE\!nd^{\mathrm{ch}}_A$ is the nonsymmetric operad with components
\[
\cE\!nd^{\mathrm{ch}}_A(n)
:=
\Hom\bigl(A^{\otimes n},\,
A((\lambda_1))\cdots((\lambda_{n-1}))\bigr),
\]
and partial compositions $f \circ_i g$ defined by operadic insertion
with the block-substitution rule of Notation~\ref{not:spectral}:
the variable $\lambda_i$ in $f$ is replaced by the block sum
$\Lambda_I = \lambda_i + \cdots + \lambda_{i+m-2}$, and the
variables of $g$ are inserted consecutively.
Associativity of the composition follows from associativity of
consecutive block-substitution. Then:
\begin{enumerate}[label=\textup{(\roman*)}]
\item An $A_\infty$ chiral algebra structure on $A$ is a degree-$+1$
 element $m = \sum_{k \ge 1} m_k \in \prod_{k \ge 1}
 \cE\!nd^{\mathrm{ch}}_A(k)$ satisfying $[m, m] = 0$, where
 the bracket is the Gerstenhaber bracket
 $[f,g] = f\{g\} - (-1)^{(|f|-1)(|g|-1)} g\{f\}$
 induced by operadic insertion.
\item The chiral Hochschild cochain complex
 $\mathrm{C}^\bullet_{\mathrm{ch}}(A, A)
 = \prod_{n \ge 0} \cE\!nd^{\mathrm{ch}}_A(n+1)[-n]$
 with differential $\delta = [m, -]$ is a brace dg algebra
 \textup{(}Volume~I, Theorem~\textup{\ref*{V1-thm:thqg-brace-dg-algebra}}\textup{)}.
\item The cohomology
 $\cZ^{\mathrm{der}}_{\mathrm{ch}}(A)
 := H^*(\mathrm{C}^\bullet_{\mathrm{ch}}(A, A), \delta)$
 is the \emph{chiral derived center} and carries a Gerstenhaber algebra
 structure. It is the universal bulk of the holomorphic-topological
 theory with boundary~$A$
 \textup{(}Volume~I, Theorem~\textup{\ref*{V1-thm:thqg-swiss-cheese}}\textup{)}.
\end{enumerate}
The reformulation aligns the axiomatic definition with the operadic
machinery of Volume~I, the open/closed factorization construction
of Part~\ref{part:swiss-cheese}.
\end{remark}

\begin{remark}[Bar--coderivation isomorphism and the
$\SCchtop$ colour assignment]
\label{rem:bar-coder-SCchtop-colour-assignment}
\index{chiral Hochschild!bar-coderivation isomorphism|textbf}%
\index{SCchtop@$\SCchtop$ colour assignment!closed = brace, open = $A$}%
The chiral Hochschild cochain complex of
\textup{Remark~\ref{rem:operadic-formulation}\,(ii)} carries an
equivalent presentation via coderivations of the bar coalgebra:
\begin{equation}\label{eq:bar-coder-iso}
 \mathrm{C}^{\bullet}_{\mathrm{ch}}(A, A)
 \;\simeq\;
 \mathrm{Coder}_0\bigl(\mathrm B^{\mathrm{ch}}(A)\bigr)
 \;\simeq\;
 \mathrm{Coder}\bigl(\barBch(A)\bigr),
 \qquad
 d_{\mathrm{Hoch}} \;=\; [D_A,\, -\,],
\end{equation}
where $\mathrm B^{\mathrm{ch}}(A)=k\oplus\barBch(A)$,
\(\barBch(A)=T_{c}(s^{-1}\bar A)\), and
\(D_A=\sum_{r\ge1}D_{m_r}\) is the degree-\(1\) structure
coderivation of the
$E_{1}$-chiral bar coalgebra of
\cref{thm:typed-boundary-holographic-realisation}\,(i), and
\(\mathrm{Coder}_0\) denotes coaugmentation-preserving coderivations.
The equivalence is the universal property of the cofree conilpotent tensor coalgebra:
a coderivation of $T_{c}(V)$ is uniquely determined by its corestriction
to $V$, which is a degree-$+1$ map $V \to V^{\otimes \le k}$ for each
$k$, i.e.\ an element of
$\prod_{n \ge 0} \cE\!nd^{\mathrm{ch}}_{A}(n + 1)[-n]$.

In the $\SCchtop$ operadic language, the chiral Deligne action of
\(\Coder_0(\mathrm B^{\mathrm{ch}}(A))\) on \(A\) is the two-colour assignment
\begin{description}
\item[\textup{closed colour}] \(\Coder_0(\mathrm B^{\mathrm{ch}}(A)) \simeq
\Zderch{A}\) acts by chiral braces, the closed colour of $\SCchtop$
(\cref{thm:chiral-higher-deligne}\,(1)--(2));
\item[\textup{open colour}] $A$ acts by its native $E_{1}$-chiral
product $m_{2}$, the open colour of $\SCchtop$
(\cref{def:e1-chiral-algebra}).
\end{description}
The mixed (open/closed) operations of $\SCchtop$ are the chiral
brace insertions \(\nu^{(m)}_{k} \colon
\Coder_0(\mathrm B^{\mathrm{ch}}(A))^{\otimes k} \otimes A^{\otimes m} \to A\)
of the proof of \cref{thm:chiral-higher-deligne}\,(2).  The precise
chiral $A_{\infty}$-Deligne theorem --- the assembly of these data
into a coherent $\SCchtop$-algebra structure --- is
\cref{thm:chd-ds-hochschild} on the chirally Koszul locus, with the
extra two-colour cobar contracting homotopy of
\cref{rem:chd-universality-hypotheses}.
\end{remark}

Definition~\ref{def:ainfty_chiral} is the coordinate expression (in spectral parameters $\lambda_i$) of a single global object on the curve~$X$: an algebra structure depending on the ordered positions of the inputs and holomorphic in their separations. The global object is named next.

\begin{definition}[$\Eone$-chiral algebra]
\label{def:e1-chiral-algebra}
An \emph{$\Eone$-chiral algebra} on a smooth curve $X$ is an
$\Eone$-algebra object in the symmetric monoidal dg category of
$\cD$-modules on~$X$, equipped with chiral operations. Concretely, it is a chiral algebra $\cA$ on~$X$
equipped with $\Ainf$ operations
\[
m_k \colon \cA^{\otimes k} \longrightarrow \cA,
\qquad k \ge 1,
\]
whose structure constants depend on \emph{ordered} configurations,
with the $\Ainf$ relations (Stasheff associahedra) holding on the
Fulton--MacPherson compactifications $\FM_k(\C)$. The operations
$m_k$ encode the OPE poles after $d\!\log$ extraction: an OPE pole
of order~$p$ contributes to~$m_p$.

This is the open colour of $\SCchtop$: the chiral factor
$\FM_k(\C)$ governs holomorphic collisions, and the $\Eone$ factor
$\Eone(m)$ records the topological ordering along~$\R$.
Definition~\ref{def:ainfty_chiral} is the coordinate expression of
this structure in spectral parameters.
\end{definition}

\begin{definition}[$\Eone$-topological algebra]
\label{def:e1-topological-algebra}
An \emph{$\Eone$-topological algebra} (equivalently, an
$\Ainf$-algebra in the classical sense) is an associative algebra up
to coherent homotopy, with no holomorphic or complex-structure
dependence. The $\Eone$ operad is modeled by ordered configurations
of intervals on~$\R$: the operation space $\Eone(k)$ is the space of
$k$~ordered disjoint subintervals of a standard interval, and the
$\Ainf$ relations arise from the cellular decomposition of the
Stasheff associahedra $K_n$.

This is the topological factor of $\SCchtop$: restricting the
two-colored operad $\SCchtop$ to operations involving only the open
colour $\mathsf{top}$, and forgetting the $\FM_k(\C)$ factor, yields
the $\Eone$ operad governing $\Eone$-topological algebras.
\end{definition}

\begin{remark}[$\Eone$-chiral vs.\ $\Eone$-topological]
\label{rem:e1-chiral-vs-topological}
The distinction is the dependence on the complex structure of~$X$.
An $\Eone$-chiral algebra depends on the complex structure: the OPE
poles are holomorphic data, and the operations $m_k$ are Laurent
polynomials in spectral parameters $\lambda_i$ encoding holomorphic
separations. An $\Eone$-topological algebra carries no such
dependence: its operations are abstract multilinear maps satisfying
the Stasheff relations without spectral parameters.

An $\Eone$-chiral algebra restricted to the topological direction
(forgetting the holomorphic $\FM_k(\C)$ factor, retaining only the
$\Eone(m)$ ordering) yields an $\Eone$-topological algebra. Operadically,
this is the forgetful functor
$\mathrm{Alg}_{\SCchtop} \to \mathrm{Alg}_{\Eone}$ discarding the
chiral colour.

A conformal vector (Virasoro element) is not part of the $\Eone$
level: neither the $\Eone$-chiral nor the $\Eone$-topological
structure carries one. The Virasoro action arises with
$\Eone$-chiral structure compatible with conformal symmetry, i.e.\
when $\cA$ is a vertex algebra with a conformal vector.
\end{remark}

The $\Eone$-chiral notion sits one step below the
$\Einf$-chiral factorization algebra, its symmetric-descent sibling.
Pole content is compatible with either structure, so the
discriminant is locality: a nonzero monodromy class on the ordered
configuration space obstructing $\Sigma_n$-descent.

\begin{proposition}[$\Eone$-chiral vs.\ $\Einf$-chiral: the locality obstruction]
\label{prop:e1-chiral-vs-e-infty-chiral-obstruction}
\leavevmode
\begin{enumerate}[label=\textup{(\roman*)}]
\item \emph{Definitional dichotomy.} An $\Einf$-chiral algebra on~$X$
is a factorization algebra on $\Ran(X)$ whose $n$-ary operations are
$\Sigma_n$-equivariant sections over the \emph{unordered} configuration
space $\Conf_n(X) = (X^n \setminus \Delta)/\Sigma_n$; an
$\Eone$-chiral algebra is a factorization algebra on the
\emph{ordered-stratified} Ran space $\Ran(X)^{\mathrm{ord}}$, whose
$n$-ary operations live on the ordered configuration space
$\Conf_n^{\mathrm{ord}}(X) = X^n \setminus \Delta$ and are
\emph{not} required to descend to the $\Sigma_n$-quotient. Poles of
the OPE are compatible with either structure; the discriminant is
\emph{locality} (symmetric descent), not pole content.
\item \emph{Operadic analogy.} $\Einf$-chiral is to $\Eone$-chiral as
commutative is to associative: an $\Einf$-chiral algebra is a
commutative algebra object in the factorization ambient
$\Fact(X)$; an $\Eone$-chiral algebra is an associative algebra
object in the ordered factorization ambient
$\Fact^{\mathrm{ord}}(X)$. The inclusion
$\Fact(X) \hookrightarrow \Fact^{\mathrm{ord}}(X)$ is the chiralisation
of the inclusion of commutative algebras into associative algebras.
\item \emph{Obstruction to promotion.} The Yangian $\Yhbar(\fg)$ with
universal $R$-matrix $R(z_1 - z_2)$ is an $\Eone$-chiral algebra that
does \emph{not} admit any $\Einf$-chiral completion. Concretely, the
two evaluation products obtained by composing along
$(z_1, z_2) \mapsto z_1 \to z_2$ and $(z_1, z_2) \mapsto z_2 \to z_1$
differ by the $R$-matrix twist:
\[
\mu(z_1, z_2) = R(z_1 - z_2) \cdot \tau \cdot \mu(z_2, z_1),
\]
and $R(z_1 - z_2) \ne R(z_2 - z_1)$ generically
(quasi-triangularity is \emph{asymmetric}, not $\Sigma_2$-equivariant).
Hence no $\Sigma_2$-equivariant lift to operations on
$\Conf_2(X)$ exists: the $\Sigma_2$-coinvariant of the ordered
operation $\mu(z_1, z_2)$ is not well-defined as a factorization
datum on the unordered stratum.
\item \emph{Factorization-algebra functor.} The functor
\[
\Phi_X : \Yhbar(\fg)\text{-}\mathrm{mod} \longrightarrow
\Fact^{\mathrm{ord}}(X)
\]
sending the Yangian to its associated $\Eone$-chiral factorization
algebra on the ordered Ran space \emph{exists}
\textup{(}Francis--Gaitsgory \cite{FG12}, IV.5, in the ordered
variant\textup{)}. The analogous functor to $\Fact(X)$ with
values in $\Einf$-chiral factorization algebras does \emph{not}
exist: the obstruction (iii) is detected by the non-vanishing of the
elementary-transposition monodromy of the local system $\cL_R$ of
Theorem~A$^{\infty,2}$, and equivalently by the fact that the
$R$-twisted $\Sigma_n$-descent of the ordered bar
$B^{\mathrm{ord}}(\Yhbar)$ does not degenerate to ordinary
$\Sigma_n$-coinvariants.
\item \emph{Contrast with the $\Einf$-chiral landscape.} The
Heisenberg, affine Kac--Moody, Virasoro, lattice, and $W$-algebra
families all carry $\Sigma_n$-equivariant operations on
$\Conf_n(X)$ with OPE poles of any order; their $R$-matrices, when
recorded, are \emph{derived} from the OPE via
$R(z) = \exp(r(z) \cdot \hbar + \cdots)$ applied to the symmetric
collision residue, hence are $\Sigma_2$-symmetric modulo the
tautological swap $\tau$. $\Einf$ admits pole-bearing realisations;
the discriminant from $\Eone$-chiral is locality (symmetric descent),
not pole content.
\end{enumerate}
\end{proposition}

\begin{remark}[Sharpness of the distinction]
\label{rem:e1-einf-chiral-sharp}
Item (iii) of Proposition~\ref{prop:e1-chiral-vs-e-infty-chiral-obstruction}
is a nonzero cohomology class — the monodromy of $\cL_R$ on
$\Conf_n^{\mathrm{ord}}(X)$ — not a presentation artifact. Completion,
localisation, or higher-categorical replacement do not kill a nonzero
local-system
monodromy. The identification
$\Eone\text{-chiral} \leftrightarrow \mathrm{Assoc}$ in
$\Fact^{\mathrm{ord}}(X)$ vs.\
$\Einf\text{-chiral} \leftrightarrow \mathrm{Comm}$ in $\Fact(X)$
is as sharp as the classical distinction between associative and
commutative algebras.
\end{remark}

\subsection{Sesquilinearity and unitality axioms}
\label{subsec:sesqui}

Definition~\ref{def:ainfty_chiral} equips $A$ with a translation
operator $\partial$ and operations $m_k$ valued in formal Laurent
series in the spectral parameters $\lambda_i$. The two data are
coupled by translation equivariance of the chiral distribution
kernels: applying $\partial$ to an input is the Fourier transform of
the infinitesimal translation acting on the corresponding collision
coordinate. In the divided-power $\lambda$-bracket normalisation this
is multiplication by the relevant spectral variable. The unit is
transparent to both structures. Sesquilinearity and unitality record
these compatibilities as identities in $\partial$ and~$\lambda_i$.

\begin{definition}[Sesquilinearity and unitality]
\label{def:sesquilinearity}
The $m_k$ satisfy the following sesquilinearity relations (generalizing the vertex-algebra relations) for $k\ge 2$:
\begin{align}
\label{eq:sesqui-left}
m_k(a_1,\dots,\partial a_i,\dots,a_k)
&= -\,\lambda_i\cdot m_k(a_1,\dots,a_i,\dots,a_k), && 1\le i\le k-1,\\
\label{eq:sesqui-right}
m_k(a_1,\dots,a_{k-1},\partial a_k)
&= \Bigl(\lambda_1+\cdots+\lambda_{k-1}+\partial\Bigr)\, m_k(a_1,\dots,a_k).
\end{align}
\noindent
\emph{Convention.}\; In equations \eqref{eq:sesqui-left}--\eqref{eq:sesqui-right} we use the
\textbf{independent-parameter convention}: $\lambda_i$ is the spectral parameter
associated to the $i$-th input (the difference $z_i - z_k$ relative to the
rightmost slot), and each $\lambda_i$ is treated as an independent formal
variable. In this convention, sesquilinearity reads
$m_k(\ldots,\partial a_i,\ldots) = -\lambda_i\cdot m_k(\ldots)$ for
$1\le i\le k-1$. The alternative \emph{consecutive-difference convention}
used in Lemma~\ref{lem:sesquilinearity_explicit} defines $\lambda_i$ as the
difference $z_i - z_{i+1}$ between adjacent insertion points.
The two conventions are related by the change of variables
\begin{equation}\label{eq:indep-to-consec}
 \lambda_i^{\mathrm{indep}} = z_i - z_k
 = \sum_{j=i}^{k-1} (z_j - z_{j+1})
 = \sum_{j=i}^{k-1} \lambda_j^{\mathrm{consec}},
 \qquad 1\le i\le k-1.
\end{equation}
In the consecutive convention, the sesquilinearity identities
\eqref{eq:sesqui-left}--\eqref{eq:sesqui-right} become: for
$1\le i\le k-1$,
$m_k(\ldots,\partial a_i,\ldots) = -(\lambda_i+\cdots+\lambda_{k-1})
\cdot m_k(\ldots)$, and for $i=k$,
$m_k(\ldots,\partial a_k) = (\partial + \sum_{j=1}^{k-1} j\,\lambda_j)
\,m_k(\ldots)$.
(The factor $j\,\lambda_j$ arises from the change of variables: $\sum_{i=1}^{k-1}\lambda_i^{\mathrm{indep}} = \sum_{i=1}^{k-1}\sum_{j=i}^{k-1}\lambda_j^{\mathrm{consec}} = \sum_{j=1}^{k-1}\bigl(\sum_{i=1}^{j}1\bigr)\lambda_j^{\mathrm{consec}} = \sum_{j=1}^{k-1}j\,\lambda_j^{\mathrm{consec}}$, where the double sum is reordered by noting that $\lambda_j^{\mathrm{consec}}$ appears for each $i\le j$.)
See Lemma~\ref{lem:sesquilinearity_explicit} for the complete derivation.

For $k=1$, $m_1=Q$ and sesquilinearity is $[Q,\partial]=0$.
There is a (strict) unit $\mathbf 1$ such that
\begin{align}
\label{eq:unit-axioms}
m_1(\mathbf 1)&=0,\qquad \partial\mathbf 1=0,\\
\label{eq:unit-m2}
m_2(\mathbf 1,a)&=a,\qquad m_2(a,\mathbf 1)=(-1)^{|a|}a,\\
\label{eq:unit-higher}
m_k(\ldots,\mathbf 1,\ldots)&=0\quad \text{for all $k\neq 2$, whenever a unit is inserted}.
\end{align}
\noindent
\emph{Sign justification.}\; The sign $(-1)^{|a|}$ in
$m_2(a,\mathbf{1})=(-1)^{|a|}a$ arises from the bar desuspension: the
operations $m_k$ are defined as the components of a degree-$+1$
coderivation on the bar complex $T^c(s^{-1}\bar A)$, and the
desuspension $s^{-1}$ produces a Koszul sign $(-1)^{|a|}$ when
passing through the first input. Concretely, the unshifted product
$\mu_2$ satisfies $\mu_2(a,\mathbf{1}) = a$, and the shifted operation
$m_2 = s \circ \mu_2 \circ (s^{-1} \otimes s^{-1})$ acquires the
sign $(-1)^{|a|}$ from the desuspension isomorphism acting on the
first tensor factor.
\end{definition}

\subsection{Planar $A_\infty$ relations with spectral substitution}
\label{subsec:ainfty-relations}

Sesquilinearity and unitality fix the interaction of each $m_k$ with $\partial$ and the unit. Closure of $\{m_k\}$ under composition remains: boundary faces of $\FM_k(\C)$ collide consecutive inputs, and Stokes' theorem forces the boundary residues to cancel. This cancellation is the Stasheff $A_\infty$ identity, upgraded by the requirement that the spectral parameter carried by the collided block is the \emph{sum} of the inner parameters (translation covariance applied to a cluster centre).

\begin{definition}[$A_\infty$ relations with spectral substitution]
\label{def:ainfty-relations}
For each $n\ge 1$, the following \emph{planar} $A_\infty$ identity holds (degree $3-n$):
\begin{multline}
\label{eq:ainfty-relation-raw}
\sum_{\substack{i+j=n+1\\ i,j\ge1}}\;\sum_{s=0}^{n-j}
(-1)^{\epsilon(s,j)}\;
m_i\!\Bigl(
a_1,\dots,a_s,\
m_j(a_{s+1},\dots,a_{s+j}\,;\,\lambda_{s+1},\dots,\lambda_{s+j-1})
\ \Big|\\
\Lambda_{[s+1,\dots,s+j]}
,\ a_{s+j+1},\dots,a_n;\ \bm\lambda^{\mathrm{out}} \Bigr)
\;=\;0.
\end{multline}
Here:
\begin{itemize}
 \item The sign $(-1)^{\epsilon(s,j)}$ is the Koszul sign for $A_\infty$ algebras in cohomological grading. The explicit formula is
 \[
 \epsilon(s,j) \;=\; |a_1|+\cdots+|a_s|,
 \]
 i.e., the parity of the total degree of the elements that the inner operation $m_j$ passes.%
 \footnote{Equation~\eqref{eq:axioms-ainfty-coderivation} is the
 suspended bar convention: the coderivation sign is
 $\sum_{r<i}(|a_r|-1)$. The expanded identity
 \eqref{eq:ainfty-relation-raw} is its desuspended coordinate form.
 The desuspension isomorphism transfers the suspended coderivation
 equation $D_A^2=0$ into the displayed Stasheff identity, with the
 signs absorbed into the coordinate operations $m_k$.
 In particular, for $j=1$: $\epsilon(s,1)=|a_1|+\cdots+|a_s|$,
 so $Q=m_1$ is a graded derivation of every~$m_k$.}
 \item The notation $m_i(-\mid \Lambda_{[s+1,\dots,s+j]},-;\bm\lambda^{\mathrm{out}})$ means:
 the $j$ consecutive inputs $a_{s+1},\dots,a_{s+j}$ are first composed by $m_j$ with the internal spectral variables $(\lambda_{s+1},\dots,\lambda_{s+j-1})$ (relative to the $a_{s+j}$ slot), resulting in one effective input placed at position $s+1$ of the outer $m_i$ with \emph{effective spectral parameter} $\Lambda_{[s+1,\dots,s+j]}$. The other outer spectral variables are given by
 \[
 \bm\lambda^{\mathrm{out}} = \bigl(\lambda_1,\dots,\lambda_s,\ \Lambda_{[s+1,\dots,s+j]},\ \lambda_{s+j+1},\dots,\lambda_{n-1}\bigr),
 \]
 i.e.\ the inner block is replaced by its sum of parameters, while the rest is unchanged.
\end{itemize}
\begin{remark}[Spectral substitution principle]
\label{rmk:spectral-substitution}
Equation \eqref{eq:ainfty-relation-raw} encodes that \emph{when a block of inputs is fused first}, the correct spectral parameter inserted in the outer operation is the \emph{sum} of the inner $\lambda$'s, reflecting translation covariance and the convention that all $\lambda_i$ are differences relative to the last slot, the algebraic avatar of momentum conservation for formal distributions.
\end{remark}
\end{definition}

\subsection{Symmetry, grading, and parity}
\label{subsec:symmetry}

The $A_\infty$ chiral algebra is a \emph{planar} (non-commutative) structure: inputs are ordered by the topological direction~$\R_t$; no symmetry under input permutation is imposed. The passage to commutative structures (PVA on cohomology, symmetric bar complex) goes through the exchange cylinder of~\S\ref{sec:PVA_descent}: in the configuration space of two points in the three-dimensional bulk $\C\times\R$, the two line orderings are joined by a path, while the holomorphic projection supplies the half-monodromy around the collision divisor.

\begin{definition}[Symmetry and grading conventions]
\label{def:symmetry-grading}
No symmetry beyond the planar order is imposed on the $m_k$ (standard for $A_\infty$ structures). The Koszul rule governs signs in all axioms above; the grading convention is cohomological.
\end{definition}

\subsection{Regular/singular decomposition of $m_2$}
\label{subsec:reg-sing}
\label{def:Laurent-projection}

The binary operation~$m_2$ is a formal Laurent series in one spectral
parameter~$\lambda$. The regular part (power series in $\lambda$)
descends to a commutative associative product on cohomology; the
singular part (negative powers of $\lambda$) descends to a
$(-1)$-shifted Lie bracket, the $\lambda$-bracket of the PVA
(Theorem~\ref{thm:cohomology_PVA}).

For $m_2(a,b)\in A((\lambda))$,
\[
m_2(a,b) = m_2^{\mathrm{reg}}(a,b)\;+\;m_2^{\mathrm{sing}}(a,b),
\]
with $m_2^{\mathrm{reg}}(a,b)\in A[\![\lambda]\!]$ and $m_2^{\mathrm{sing}}(a,b)\in \lambda^{-1}A[\![\lambda^{-1}]\!]$.

\subsection{Consequences: product and bracket on cohomology}
\label{subsec:product-bracket}
\label{def:product}
\label{def:lambda-bracket}

The regular and singular parts of $m_2$ carry distinct algebraic
content: the regular part descends to a graded-commutative
associative product on cohomology; the singular part descends to a
$(-1)$-shifted Lie bracket, the $\lambda$-bracket of vertex-algebra
theory. The two descents are verified in
Theorem~\textup{\ref{thm:cohomology_PVA}}.

For $[a],[b]\in H^\bullet(A,Q)$:
\begin{align}
\label{eq:prod-def}
[a]\cdot[b] &:= \Bigl[\, \frac12\bigl(m_2^{\mathrm{reg}}(a,b) + (-1)^{|a||b|}m_2^{\mathrm{reg}}(b,a)\bigr)\,\Bigr],\\
\label{eq:bracket-def}
\{[a]{}_\lambda [b]\} &:= \Bigl[\, \frac12\bigl(m_2^{\mathrm{sing}}(a,b) - (-1)^{|a||b|}\; \mathrm{e}^{-(2\lambda+\partial)\partial_\lambda}\,m_2^{\mathrm{sing}}(b,a)\bigr)\,\Bigr].
\end{align}
The exponential operator encodes the transformation $\lambda\mapsto -\lambda-\partial$ enforced by vertex skewsymmetry. (Action on formal series: $\mathrm{e}^{\alpha\partial_\lambda}f(\lambda)=f(\lambda+\alpha)$.)

\smallskip\noindent
\emph{The exponential operator.}\;
$\mathrm{e}^{\alpha\partial_\lambda}$ acts on $f(\lambda)$ by
$f(\lambda+\alpha)$. Setting $\alpha = -(2\lambda+\partial)$ gives
$f(\lambda) \mapsto f(-\lambda-\partial)$, the vertex-algebra
skew-symmetry substitution. The visible exchange sign is the
ordinary Koszul sign $(-1)^{|a||b|}$; shifted signs enter the
Jacobi and Leibniz identities through the bar desuspension.

\smallskip\noindent
\emph{Symmetrisation in the product.}\;
$m_2$ is not graded-commutative — it encodes the full OPE, not
just the commutative normally-ordered product — so the
symmetrisation in~\eqref{eq:prod-def} is required. The factor
$\frac{1}{2}$ is well-defined in characteristic~$0$ (a standing
hypothesis on the ground field, cf.\
Definition~\ref{def:ainfty_chiral}). If $m_2^{\mathrm{reg}}$ is
already graded-commutative (Heisenberg algebra), the symmetrisation
is redundant and $[a]\cdot[b] = [m_2^{\mathrm{reg}}(a,b)]$.

The PVA axioms on $H^\bullet(A, Q)$ are verified in
Theorem~\ref{thm:cohomology_PVA}.

\subsection{The QME as the $A_\infty$ identity}
\label{subsec:qme-ainfty}

The $A_\infty$ identities of Definition~\ref{def:ainfty-relations} were derived above as the Stasheff closure forced by Stokes' theorem on $\FM_k(\C)$. For an $A_\infty$ chiral algebra arising from a BV quantization of a 3d holomorphic-topological field theory, the same identities have a second life as the Feynman-diagram expansion of the quantum master equation (QME). The coincidence is not a coincidence: each $A_\infty$ relation is the QME at a fixed diagrammatic weight.

\begin{lemma}[QME implies $A_\infty$ relations]
\label{lem:QME_to_Ainfty}
For a physical realisation satisfying Theorem~\ref{thm:physics-bridge}:
The quantum master equation $\frac{1}{2}\{S_{\mathrm{eff}},S_{\mathrm{eff}}\}_{\mathrm{BV}} + \hbar\Delta S_{\mathrm{eff}}=0$, expanded in Feynman diagrams, is equivalent to the $A_\infty$ identities \eqref{eq:ainfty-relation-raw} for the operations $m_k$ extracted from $S_{\mathrm{eff}}$.
\end{lemma}

\begin{proof}
Expand $S_{\mathrm{eff}} = \sum_k \frac{1}{k!} S_k$ where $S_k$ encodes $k$-ary operations. The full QME $\tfrac{1}{2}\{S_{\mathrm{eff}},S_{\mathrm{eff}}\}_{\mathrm{BV}} + \hbar\Delta S_{\mathrm{eff}} = 0$ at each degree $n$ yields precisely the sum over compositions $m_i \circ m_j$ with $i+j=n+1$, matching \eqref{eq:ainfty-relation-raw}. The complete argument appears in \cite{CG17}*{\S5.4} and \cite{CDG20}*{\S3}.
\end{proof}

\begin{proposition}[Definition~\ref{def:ainfty_chiral} captures the physical structure]
\label{prop:physical_justification}
For physical realisations satisfying Theorem~\ref{thm:physics-bridge}, Definition~\ref{def:ainfty_chiral} formalizes the $A_\infty$ chiral algebra structure of Costello--Gwilliam--Gaiotto--Dimofte:

\begin{enumerate}[label=(\roman*)]
\item The operations $m_k$ encode the factorization product on $\Conf_k(\C)$, recovering the Costello--Gwilliam structure \cite{CG17}.

\item The $A_\infty$ identity \eqref{eq:ainfty-relation-raw} is the quantum master equation
$\frac{1}{2}\{S_{\mathrm{eff}}, S_{\mathrm{eff}}\}_{\mathrm{BV}} + \hbar \Delta S_{\mathrm{eff}} = 0$
expanded in Feynman diagrams.

\item The $\lambda$-bracket and form integrands match \cite{CDG20} and Khan--Zeng \cite{KZ25}: spectral parameters $\lambda_i$ encode derivatives via $e^{\lambda z}$ generating functions.

\item The definition recovers known structures:
binary descent gives ordinary PVAs (Theorem~\ref{thm:cohomology_PVA});
the additional condition $m_{k \geq 3}=0$ gives formal comparison
models (Example~\ref{rem:free_chiral});
one-loop-finite holomorphic theories give Williams structures \cite{GGW21};
cohomology $H^\bullet(\A,Q)$ yields PVAs matching \cite{CDG20} (Theorem~\ref{thm:cohomology_PVA}).

\item The definition generalizes Tamarkin's chiral Hochschild complex \cite{Tam02}: it is the 3d analogue of 2d Hochschild theory from Poisson sigma models (Section~\ref{sec:chiral_hochschild_expanded}).
\end{enumerate}
\end{proposition}

\begin{remark}[Comparison with the physical source]
\label{rem:physical-source-comparison}
Costello--Dimofte--Gaiotto \cite{CDG20} isolate the need for an
$A_\infty$ analogue of a vertex algebra in holomorphic-topological
field theory. Definition~\ref{def:ainfty_chiral} gives the Vol~II
model: spectral $A_\infty$ operations, sesquilinearity, unit,
Stasheff identities, and the logarithmic $\SCchtop$ factorisation
needed for PVA descent.
\end{remark}

\subsection{Factorization axiom and cluster limits}

As insertion points collide, the operations decompose according to the boundary stratification of $\FM_k(\C)$. This \emph{cluster factorization} axiom is the bridge between the $m_k$ and the operadic structure of~$\OHT$.

\begin{theorem}[Cluster factorization in homotopy-commutative form]
\label{thm:cluster_factorization_homotopy}
Let $(\A, \{m_k\})$ be a logarithmic $\SCchtop$-algebra (a $C_\ast(W(\SCchtop))$-algebra whose weight forms factor as $\omega_k = \omega_k^{\mathrm{hol}} \otimes \omega_k^{\mathrm{top}}$ on $\FM_k(\C) \times \Conf_k(\R)$; the full definition is Definition~\ref{def:log-SC-algebra}) and let $\pi = \{S_1,\ldots,S_r\}$ be a partition of $\{1,\ldots,k\}$ with $|S_j| = k_j$. As points $\{z_i\}_{i \in S_j}$ collide to cluster centers $w_j$ (a limit within $\FM_k(\C)$), the operations satisfy
\begin{equation}
\label{eq:cluster_factorization}
\lim_{z_i \to w_j} m_k(O_1,\ldots,O_k) = \sum_{\sigma \in \Sh(\pi)} (-1)^{\varepsilon(\sigma;O)} \, m_r\big(O'_{\sigma(1)},\ldots,O'_{\sigma(r)}\big),
\end{equation}
where $O'_j := m_{k_j}(O_{i_1},\ldots,O_{i_{k_j}})$ for $S_j = \{i_1 < \cdots < i_{k_j}\}$, and spectral parameters compose via
\begin{equation}
\label{eq:spectral_composition}
\lambda_{\text{cluster},j} = \sum_{i \in S_j} \lambda_i + \partial_{w_j}.
\end{equation}
Here $\lambda_i$ for $i \in S_j$ denotes the inner spectral parameters relative to the cluster center; the sum runs over $|S_j| - 1$ independent parameters (the last being determined by the constraint $\sum_{i \in S_j} \lambda_i = \Lambda_j - \partial_{w_j}$).
The shuffle sum $\sum_{\sigma\in\Sh(\pi)}$ arises because the Swiss-cheese
operations carry a homotopy-commutative structure induced by the symmetric
group action on $\FM_k(\C)$. For the planar (non-commutative) specialization,
only the identity shuffle $\sigma=\mathrm{id}$ contributes, and the sum
reduces to a single term.
\end{theorem}

\begin{proof}
The boundary factorisation has geometric, analytic, and spectral
components.

\textbf{Boundary stratum.} The boundary stratum $D_\pi \subset \partial \FM_k(\C)$ corresponding to partition $\pi$ is canonically isomorphic to
\[
D_\pi \cong \FM_r(\C) \times \prod_{j=1}^r
\FM_{k_j}^{\mathrm{red}}(\C),
\]
recording cluster centers $\{w_j\}$ and relative positions within each cluster.

\textbf{Logarithmic residue.} Let $S\subset\{1,\ldots,k\}$ be one
collision block and choose the inward radial coordinate
$\varepsilon_S\geq0$ on the corresponding boundary divisor
$D_S$.  In the FM logarithmic chart the weight form has the local
normal form
\[
\omega_k=d\log\varepsilon_S\wedge\beta_S+\alpha_S,
\qquad
\alpha_S,\beta_S\ \text{smooth at }\varepsilon_S=0,
\]
and the residue used in the cluster operation is
\begin{equation}
\label{eq:cluster-fm-residue}
\Res_{D_S}(\omega_k)=\beta_S\big|_{\varepsilon_S=0}
=
\lim_{\varepsilon_S\to0}
\bigl(\varepsilon_S\iota_{\partial_{\varepsilon_S}}\omega_k\bigr)
\big|_{D_S}.
\end{equation}
The naive contraction
$\iota_{\partial_{\varepsilon_S}}\omega_k|_{D_S}$ is not used: for
logarithmic forms it has a pole of order $\varepsilon_S^{-1}$.
The boundary orientation is
\begin{equation}
\label{eq:cluster-fm-orientation}
\operatorname{or}\bigl(\FM_k(\C)\bigr)
=(-d\varepsilon_S)\wedge\operatorname{or}(D_S).
\end{equation}
Thus the incidence sign $\epsilon_{D_S}$ in Stokes' theorem is fixed
by moving the logarithmic normal factor to the outward-normal slot.

For the full partition $\pi$ the ordered iterated residue along the
normal-crossing stratum $D_\pi$ factors as
\[
\Res_{D_\pi}(\omega_k)
=
(-1)^{\epsilon_\pi}
\omega_r(O'_1,\ldots,O'_r) \wedge
\bigwedge_{j=1}^r \omega_{k_j}(\{O_i\}_{i \in S_j}),
\]
where $\epsilon_\pi$ is the incidence sign determined by
\eqref{eq:cluster-fm-orientation} and the Koszul permutation of
operator labels.  Equivalently, the collision part of the bar
coderivation is
\begin{equation}
\label{eq:cluster-collision-differential}
D_{\mathrm{coll}}
=
\sum_{\substack{S\subset\{1,\ldots,k\}\\ |S|\ge2}}
\epsilon_{D_S}\,\Res_{D_S}\otimes m_S .
\end{equation}
The identity $D_{\mathrm{coll}}^2=0$ is the image of
$\partial^2[\FM_k(\C)]=0$.  At a codimension-two corner
$D_S\cap D_T$, the two ordered normal approaches contribute
\[
\int_{(D_S)|_{D_T}}\Res_{D_T}\Res_{D_S}(\omega_k)
+
\int_{(D_T)|_{D_S}}\Res_{D_S}\Res_{D_T}(\omega_k)=0,
\]
because the normal covectors anticommute,
$d\varepsilon_T\wedge d\varepsilon_S
=-d\varepsilon_S\wedge d\varepsilon_T$.

\textbf{Spectral composition.} Holomorphic generating functions $e^{\sum_i \lambda_i z_i}$ degenerate to $e^{(\sum_{i \in S_j} \lambda_i) w_j}$ times relative terms, yielding (\ref{eq:spectral_composition}).
\end{proof}

\begin{remark}[Nonabelian Hodge shadow problem]
\label{rem:nonabelian_hodge_shadow}
The shadow obstruction tower of the chiral factor defines the following
nonabelian Hodge enhancement problem: construct a triple
$(\Theta_\cA,M_\cA,h_\cA)$ consisting of the MC element
$\Theta_\cA$ on the flat side, a shadow Higgs field $M_\cA$, and a
Shapovalov-type harmonic form $h_\cA$, such that
$d=1+d_{\mathrm{arith}}+d_{\mathrm{alg}}$ is the Hodge decomposition of
the triple. Cluster factorization does not use this enhancement; its
construction belongs to the nonabelian Hodge package of Volume~I.
\end{remark}

\begin{remark}[Kac--Kazhdan / Shapovalov determinant: scope and usage]
\label{rem:shapovalov-determinant-scope}
\index{Kac-Kazhdan determinant}\index{Shapovalov determinant!scope}%
The Kac--Kazhdan / Shapovalov determinant appears as a
\emph{vanishing-locus control} underlying the class-M shadow tower,
the critical-level (class FF) structure, and the minimal-model
null-vector corrections. The following three items fix the scope.

\smallskip
\textbf{(a) Classical Shapovalov (Virasoro).} On the weight-$n$ graded piece
of the Verma module $M(c,h)$, the Shapovalov form $G_n(c,h)$ has determinant
\[
\det G_n(c,h) \;=\; \prod_{\substack{p,q\ge 1\\ pq\le n}}
 \bigl(h - h_{p,q}(c)\bigr)^{p(n-pq)},
\qquad
h_{p,q}(c) \;=\; \tfrac{1}{4}\bigl(\alpha_+ p + \alpha_- q\bigr)^2 - \tfrac{(\alpha_+ + \alpha_-)^2}{16},
\]
with $\alpha_\pm$ determined by $c = 13 - 6(\alpha_+^2 + \alpha_-^2)$
(Kac 1978; Feigin--Fuchs determinant formula).
The affine-Lie-algebra generalisation is Kac--Kazhdan 1979
\cite{KacKazhdan79}: on a Verma module $M(\lambda)$ of $\widehat\fg$,
\(
\det G_{\lambda - \mu} \;=\;
\prod_{\alpha \in \widehat\Delta_+^{\mathrm{re}}}\prod_{n\ge 1}
 \bigl(\langle \lambda+\rho, \alpha\rangle - n\alpha(\alpha)/2\bigr)^{\mathrm{mult}}
\times
\prod_{\delta}(\langle\lambda+\rho,\delta\rangle)^{\dim \fg \cdot P(\cdot)},
\)
the second factor the imaginary-root contribution absent from the finite case.

\smallskip
\textbf{(b) Shadow-tower consequence (class~M, $r_{\max}=\infty$).}
The class-M assignment for Vir$_c$ at generic $c$ and for
$\mathcal W_N$ at generic $c$ rests on the following structural fact:
the Kac--Kazhdan vanishing locus
$\mathcal Z_{KK} = \{(c,h) : \det G_n(c,h) = 0\text{ for some }n\}$
is a countable union of \emph{positive-codimension} analytic subvarieties
of $(c,h)$-parameter space. Away from $\mathcal Z_{KK}$ the Verma module
is irreducible, no singular vectors appear, and the shadow-tower
coefficients $S_r$ (Vol.~I, Theorem~D) do not terminate: this is the
$r_{\max}=\infty$ signature. On $\mathcal Z_{KK}$ (the minimal-model / admissible
locus), the simple quotient acquires null vectors; the shadow tower is
corrected \emph{not by termination} but by the explicit null-vector ideal
(Arakawa $C_2$-cofiniteness for the simple quotient;
Proposition~\ref*{V1-prop:minimal-model-non-koszul}).
Class M thus persists on the simple quotient at minimal-model loci.

\smallskip
\textbf{(c) $\mathcal W$-algebra analog.} For $\mathcal W_N = \mathcal W(\sln_N)$
at principal nilpotent, the Shapovalov-type determinant on the $N-1$ highest-weight parameters is
given by the finite-$\mathcal W$ and quantum-$\mathcal W$ comparison
papers of de~Boer--Tjin~\cite{deBoerTjin93,deBoerTjinRelation93};
Bouwknegt--Schoutens~\cite{BMP-Wsymmetry} review the
Drinfeld--Sokolov and coset constructions of the same family.
The vanishing locus generalises $h_{p,q}(c)$ to a Weyl-group orbit on
$\mathfrak h^\ast$; positive-codimension away from BRST-degenerate walls.
This confirms the class-M assignment for $\mathcal W_N$ at generic central
charge and identifies the minimal-$\mathcal W$-model walls as the
$\mathcal W$-analog of the $(p,q)$-lattice.

\smallskip
\textbf{Inputs.} Item~(a) is the classical determinant formula of
Kac~1978 and Kac--Kazhdan~1979; item~(b) combines the Vol.~I class-M
calculation with Kac--Kazhdan vanishing; item~(c) is the
$\mathcal W_N$ analogue of de~Boer--Tjin and Bouwknegt--Schoutens.
The remark consolidates the Kac--Kazhdan/Shapovalov input used
downstream in this volume, and fixes the positive-codimension
convention invoked (implicitly) by the $r_{\max}=\infty$ definition
of class~M.
\end{remark}

\begin{remark}[Factorization as operadic composition]
\label{rem:factorization-as-operadic-composition}
Cluster factorization~\eqref{eq:cluster_factorization} is the statement that $\A$ is an algebra over the chain operad $\OHT = C_\ast(W(\SCchtop))$: boundary faces in $\FM_k(\C)$ encode operadic compositions in the closed colour, ordered gluings along~$\R$ encode open-colour compositions, and Stokes' theorem (Section~\ref{sec:FM_calculus}) supplies the coherence.
\end{remark}

\subsection{Explicit coend formula and operadic composition}

The operadic definition of $\OHT$ has a concrete incarnation as a coend producing the operations~$m_k$.

\begin{construction}[Coend formula for $m_k$]
\label{const:coend-formula-mk}
Using the $W$--model, the universal element determining $m_k$ is
\[
 m_k = \int^{T \in \mathrm{Trees}_k} \!\!\! C_\ast\!\bigl(W(\SCchtop)(T)\bigr)\otimes_{\mathrm{Aut}(T)}
 \bigotimes_{v \in \mathrm{Vert}(T)} \Bigl(
 C_\ast\!\bigl(\FM_{\mathrm{in}(v)}(\C)\bigr)\,\widehat{\otimes}\,
 C_\ast\!\bigl(\Conf_{\mathrm{out}(v)}(\R)\bigr)
 \Bigr)\, .
\]
Closed-colour factors record holomorphic collisions; open-colour factors record ordered
times; the $W$--parameters on internal edges control homotopies.
The coend exists because the source category $\mathrm{Trees}_k$ is finite
(finitely many labelled rooted trees with $k$ leaves) and all tensor
products are over a field of characteristic~$0$.
Remark~\ref{rem:bv-resolution-details} clarifies: the vertex-wise factors in the formula are
determined by the output colour of each vertex, not by an independent
product decomposition.

\begin{example}[Binary Operation $m_2$ from the Coend]
For $k=2$, the only tree is the corolla with one internal vertex and two leaves. The coend formula yields
\[
m_2 = \cP^{\mathrm{ch}}(2) \otimes \cE_1(1),
\]
where $\cP^{\mathrm{ch}}(2)$ encodes the binary chiral operation (holomorphic collision of two points) and $\cE_1(1)$ provides the topological ordering (represented by a single edge/interval).

Explicitly, for operators $a, b \in \A$ at positions $(t_1, z_1), (t_2, z_2)$:
\[
m_2(a,b) = \int_{\Gamma \subset \FM_2(\C)} \int_{t_1 < t_2} K(z_1 - z_2, t_1 - t_2) \, a(t_1,z_1) \otimes b(t_2,z_2) \, dt_1 \, dt_2,
\]
where $K$ is the propagator kernel; for physical realizations it is supplied by Theorem~\ref{thm:physics-bridge}.
This integral formula applies to physical realizations satisfying Theorem~\ref{thm:physics-bridge}; for abstract logarithmic $\SCchtop$-algebras, the operation $m_2$ is defined directly by the operad action.
\end{example}

\begin{remark}[Boardman--Vogt resolution details]
\label{rem:bv-resolution-details}
Each internal edge $e$ of a tree $T$ carries a parameter $s_e\in[0,1]$ encoding the time of
composition. In the two-colour model one works with $W(\SCchtop)(T)$; there is no meaningful
vertex-wise factor $(P_{\mathrm{ch}}\rtimes E_1)$. Vertex decorations are determined by the output
colour: closed vertices carry $\FM(\C)$-data; open vertices carry $\FM(\C)\times \Conf(\R)$-data. The
boundary identifications implement coherent sequential/simultaneous compositions.
\end{remark}
\end{construction}

\subsection{Sesquilinearity: Complete Proof}

The sesquilinearity identities
\eqref{eq:sesqui-left}--\eqref{eq:sesqui-right} were recorded above
as axioms; the following lemma derives them from one principle:
$\C$-translation equivariance of the chiral distribution kernels.
The proof is performed on the Fourier/divided-power coefficients.
There is no intermediate identity with
\(\partial_{\lambda_i}m_k\); the translation operator acts on the
distributional mode and Fourier transforms to multiplication by the
spectral variable. A convention-reconciliation paragraph at the end
bridges the independent-parameter form
\eqref{eq:sesqui-left}--\eqref{eq:sesqui-right} and the
consecutive-difference form
\eqref{eq:sesqui_explicit_1}--\eqref{eq:sesqui_explicit_k}.

\begin{lemma}[Sesquilinearity from translation equivariance]
\label{lem:sesquilinearity_explicit}
Here we state the result in the \textbf{consecutive-difference
convention}: the spectral parameter $\lambda_i$ is associated to
$z_i-z_{i+1}$.  This convention is related to the
independent-parameter convention of
equations~\eqref{eq:sesqui-left}--\eqref{eq:sesqui-right} by the
linear change of variables~\eqref{eq:indep-to-consec}.

Let $\partial : \A \to \A$ denote holomorphic translation in $z$. Then the operations $m_k$ satisfy:
\begin{align}
\label{eq:sesqui_explicit_1}
m_k(\partial a_1, a_2, \ldots, a_k)
 &= -(\lambda_1 + \lambda_2 + \cdots + \lambda_{k-1})\, m_k(a_1, a_2, \ldots, a_k),\\
\label{eq:sesqui_explicit_i}
m_k(a_1, \ldots, a_{i-1}, \partial a_i, a_{i+1}, \ldots, a_k)
 &= -(\lambda_i + \lambda_{i+1} + \cdots + \lambda_{k-1})\, m_k(a_1, \ldots, a_k),
 \quad 2\le i\le k-1,\\
\label{eq:sesqui_explicit_k}
m_k(a_1, \ldots, a_{k-1}, \partial a_k)
 &= \Bigl(\partial + \sum_{j=1}^{k-1} j\,\lambda_j\Bigr)\, m_k(a_1, \ldots, a_k).
\end{align}
Note that \eqref{eq:sesqui_explicit_1} is the $i=1$ case of \eqref{eq:sesqui_explicit_i};
for $k=2$ it reduces to $-\lambda_1\,m_2$.
These are obtained from \eqref{eq:sesqui-left}--\eqref{eq:sesqui-right}
via the change of variables~\eqref{eq:indep-to-consec}.
\end{lemma}

\begin{proof}
Use the independent variables
\(\mu_i=z_i-z_k\), \(1\le i\le k-1\), first.  In the
Fourier/divided-power convention for chiral operations,
\[
m_k(a_1,\ldots,a_k;\mu)
=
\sum_{\mathbf n\in\mathbb N^{k-1}}
m_{k,\mathbf n}(a_1,\ldots,a_k)
\prod_{r=1}^{k-1}\frac{\mu_r^{n_r}}{n_r!}.
\]
The coefficient \(m_{k,\mathbf n}\) is the operation paired with the
normalised collision mode of multi-index \(\mathbf n\). Translation
equivariance of chiral distributions gives the multivariable
divided-power mode identity
\[
m_{k,\mathbf n}(a_1,\ldots,\partial a_i,\ldots,a_k)
=
-\,n_i\,m_{k,\mathbf n-\mathbf e_i}(a_1,\ldots,a_k),
\qquad 1\le i<k,
\]
where the right side is zero when \(n_i=0\).  This is the
multi-input form of the vertex-algebra covariance
\((T a)_{(n)}=-n\,a_{(n-1)}\), obtained by applying the
infinitesimal translation generator to the input and integrating by
parts on the collision coordinate. Multiplying by the divided powers
and summing gives
\[
m_k(a_1,\ldots,\partial a_i,\ldots,a_k;\mu)
=-\mu_i\,m_k(a_1,\ldots,a_k;\mu),
\qquad 1\le i<k.
\]
For the last input, \(\partial_{z_k}\) differentiates the output
section as well as the independent variables
\(\mu_r=z_r-z_k\). The same translation-equivariance identity
therefore has one additional output term:
\[
m_{k,\mathbf n}(a_1,\ldots,a_{k-1},\partial a_k)
=
\partial_A m_{k,\mathbf n}(a_1,\ldots,a_k)
\;+\;
\sum_{r=1}^{k-1} n_r\,
m_{k,\mathbf n-\mathbf e_r}(a_1,\ldots,a_k).
\]
After summing the divided-power series this becomes
\[
m_k(a_1,\ldots,a_{k-1},\partial a_k;\mu)
=
\left(\partial_A+\sum_{r=1}^{k-1}\mu_r\right)
m_k(a_1,\ldots,a_k;\mu),
\]
which is the independent-parameter form
\eqref{eq:sesqui-left}--\eqref{eq:sesqui-right}.

It remains only to write the same formula in consecutive variables.
Since
\[
\mu_i=\lambda_i+\lambda_{i+1}+\cdots+\lambda_{k-1},
\]
the identities for \(1\le i<k\) become
\[
m_k(\ldots,\partial a_i,\ldots)
=-(\lambda_i+\cdots+\lambda_{k-1})m_k(\ldots,a_i,\ldots),
\]
which is \eqref{eq:sesqui_explicit_1} for \(i=1\) and
\eqref{eq:sesqui_explicit_i} for \(2\le i\le k-1\).  For the last
input,
\[
\sum_{i=1}^{k-1}\mu_i
=
\sum_{i=1}^{k-1}\sum_{j=i}^{k-1}\lambda_j
=
\sum_{j=1}^{k-1} j\,\lambda_j,
\]
and this gives \eqref{eq:sesqui_explicit_k}.  No formal derivative
\(\partial_{\lambda_i}\) appears in the \(\lambda\)-bracket
normalisation; the operator \(\partial_A\) acts on coefficients and
the variables \(\lambda_i\) act by multiplication in the divided-power
generating series.
\end{proof}
\subsection{From $W(\mathsf{SC}^{\mathrm{ch,top}})$ to axioms and back}
\label{sec:operad-to-axioms}
\label{sec:equivalence}

\subsubsection{From the operad to explicit $m_k$ with spectral substitution}
\label{subsec:operad-to-mk}
Let $(A_{\mathsf{ch}},A_{\mathsf{top}})$ be a $C_\ast\!\bigl(W(\mathsf{SC}^{\mathrm{ch,top}})\bigr)$--algebra with underlying complex $A$ and translation operator $\partial$ (from the holomorphic \(\C\)--action). Choose, for each $k\ge1$, a fundamental chain $[\FM_k(\C)]\in C_\ast(\FM_k(\C))$ and similarly for $E_1$ (recording the $t$--ordering). Define
\[
m_k(a_1,\dots,a_k)
:= \bigl\langle\ \rho\bigl([\FM_k(\C)]\times [E_1(\text{ordering})]\bigr)\,,\, a_1\otimes\cdots\otimes a_k\ \bigr\rangle
\ \in\ A((\lambda_1))\cdots((\lambda_{k-1})),
\]
where $\rho$ is the operad action and the \(\lambda\)--dependence is the generating function of the log--forms (residue degrees) on $\FM_k(\C)$ prescribed by the operad.
Then:
\begin{enumerate}[label=(\alph*)]
 \item sesquilinearity \eqref{eq:sesqui-left}--\eqref{eq:sesqui-right}
 follows from equivariance under translations in $\C$: on divided-power
 modes it is
 \(m_{k,\mathbf n}(\ldots,\partial a_i,\ldots)
 =-n_i m_{k,\mathbf n-\mathbf e_i}(\ldots)\), and summing the
 generating series gives multiplication by the corresponding
 spectral variable,
 \item the spectral substitution law in \eqref{eq:ainfty-relation-raw} is the image of the boundary identification of faces in $\partial\FM_k(\C)$ where a consecutive block collapses (the effective parameter is the sum of inner ones).
\end{enumerate}

\begin{proposition}[Operad $\Rightarrow$ axioms]
\label{prop:operad-implies-axioms}
The $m_k$ produced by the $W(\mathsf{SC}^{\mathrm{ch,top}})$--action satisfy the axioms in Section~\ref{sec:axioms-Ainfty-chiral}.
\end{proposition}

\begin{proof}
Let $(A_{\mathsf{ch}}, A_{\mathsf{top}})$ be a $C_\ast(W(\mathsf{SC}^{\mathrm{ch,top}}))$-algebra with operad action $\rho$, and let $m_k$ be defined by evaluation on fundamental chains as in \S\ref{subsec:operad-to-mk}:
\[
m_k(a_1,\ldots,a_k) := \bigl\langle \rho\bigl([\FM_k(\C)] \times [E_1(\mathrm{ord})]\bigr),\, a_1 \otimes \cdots \otimes a_k \bigr\rangle.
\]
The operad action gives each axiom of
Definition~\ref{def:ainfty_chiral}.

\medskip
\textbf{(I) Unitality.}
The operadic unit in $\mathsf{SC}^{\mathrm{ch,top}}$ is the identity element $\mathbf{1} \in \mathsf{SC}^{\mathrm{ch,top}}(1)$, represented by the point $\FM_1(\C) = \mathrm{pt}$. By the operadic unit axiom, $\rho(\mathbf{1}) = \mathrm{id}_A$, so $m_1(\mathbf{1}) = Q(\mathbf{1}) = 0$ (since the differential on the operad restricts to zero on the unit). For $m_2$: the operadic composition $\mathbf{1} \circ_1 c_2 = c_2 = c_2 \circ_2 \mathbf{1}$ (inserting the unit into either slot of the binary corolla) gives $m_2(\mathbf{1}, a) = a$ and $m_2(a, \mathbf{1}) = (-1)^{|a|} a$, where the sign arises from the Koszul rule for the degree-$(-1)$ operation $m_2$ passing element $a$. For $k \geq 3$: the operadic composition $c_k \circ_i \mathbf{1}$ factors through $c_{k-1}$ by the operadic unit law, but the corollas $c_k$ for distinct $k$ are independent generators of $W(\mathsf{SC}^{\mathrm{ch,top}})$, and the unit insertion collapses $c_k \circ_i \mathbf{1}$ to a face of the $W$-construction with Boardman--Vogt parameter $s = 0$, which is $\partial$-exact. Hence $m_k(\ldots, \mathbf{1}, \ldots) = 0$ for $k \geq 3$. This establishes~\eqref{eq:unit-axioms}--\eqref{eq:unit-higher}.

\medskip
\textbf{(II) Sesquilinearity from translation equivariance.}
The group \(\C\) acts on the chiral distribution kernels by
translation, and the infinitesimal action on \(A\) is \(\partial\).
Writing the operation in independent variables \(\mu_i=z_i-z_k\), the
Fourier/divided-power coefficients satisfy
\[
m_{k,\mathbf n}(\ldots,\partial a_i,\ldots)
=-n_i\,m_{k,\mathbf n-\mathbf e_i}(\ldots),
\quad i<k,
\]
and
\[
m_{k,\mathbf n}(\ldots,\partial a_k)
=\partial_A m_{k,\mathbf n}(\ldots)
\;+\;
\sum_i n_i\,m_{k,\mathbf n-\mathbf e_i}(\ldots).
\]
Summing the divided-power series gives
\eqref{eq:sesqui-left}--\eqref{eq:sesqui-right}.  Passing to
consecutive variables is the linear substitution
\(\mu_i=\lambda_i+\cdots+\lambda_{k-1}\), hence
\(\sum_i\mu_i=\sum_{j=1}^{k-1}j\lambda_j\), exactly as in
Lemma~\ref{lem:sesquilinearity_explicit}.

\medskip
\textbf{(III) $A_\infty$ relations and spectral substitution from boundary faces.}
The differential on $C_\ast(W(\mathsf{SC}^{\mathrm{ch,top}}))$ is the boundary operator $\partial$ on chains of the Boardman--Vogt resolution. On the fundamental chain $[\FM_k(\C)]$, the boundary is $\partial[\FM_k(\C)] = \sum_S \epsilon_{D_S}\,[D_S]$, summing over collision divisors $D_S$ with $|S| \geq 2$. Since $\rho$ is a chain map ($\rho \circ \partial = d \circ \rho$, where $d$ on $\End(A)$ is the commutator with $Q$), applying $\rho$ to both sides gives:
\[
d\bigl(\rho([\FM_k(\C)])\bigr) = \sum_S \epsilon_{D_S}\, \rho([D_S]).
\]
Each face $D_S \cong \FM_{k-|S|+1}(\C) \times \FM_{|S|}^{\mathrm{red}}(\C)$ factorizes $\rho$ by operadic composition:
\[
\rho([D_S]) = \rho([\FM_{k-|S|+1}(\C)]) \circ_{\mathrm{ins}} \rho([\FM_{|S|}(\C)]) = m_{k-|S|+1}(\ldots, m_{|S|}(\ldots), \ldots).
\]
The identity $d^2 = 0$ on $\End(A)$ (equivalently $\partial^2 = 0$ on chains) then yields
\[
0 = d^2\bigl(\rho([\FM_k(\C)])\bigr) = \sum_{i+j=k+1}\sum_{s=0}^{k-j} (-1)^{\epsilon(s,j)}\, m_i(\ldots, m_j(\ldots), \ldots),
\]
which is the $A_\infty$ identity~\eqref{eq:ainfty-relation-raw}. Only
consecutive subsets $S$ contribute, by
Theorem~\ref{thm:stokes_arnold}.

The \emph{spectral substitution} law follows from the geometry of the boundary identification. On $D_S$ for $S = \{s{+}1,\ldots,s{+}\ell\}$, the cluster center $u_S$ replaces the $\ell$ points of $S$ by a single effective point. The relative coordinate from $u_S$ to the next external point is $u_S - z_{s+\ell} = \sum_{p=s+1}^{s+\ell-1}(z_p - z_{p+1}) = \Lambda_S$, so the outer operation receives effective spectral parameter $\Lambda_S = \lambda_{s+1} + \cdots + \lambda_{s+\ell-1}$, while the remaining outer parameters are unchanged. This matches the spectral substitution in~\eqref{eq:ainfty-relation-raw} and Remark~\ref{rmk:spectral-substitution}.

\medskip
\textbf{(IV) Degree.}
The degree of $m_k$ is $2 - k$. The bar complex $\barB(\cA) = T^c(s^{-1}\bar{\cA})$ carries a total differential $d_{\barB} = \sum_k d_k$ of cohomological degree $+1$. Each $d_k$ is the coderivation with projection
$d_k|_{(s^{-1}\bar{\cA})^{\otimes k}} = s^{-1} \circ m_k \circ s^{\otimes k}\colon (s^{-1}\cA)^{\otimes k} \to s^{-1}\cA$.
The degree of this composition is $\deg(s^{-1}) + \deg(m_k) + k\cdot\deg(s) = -1 + \deg(m_k) + k$, and requiring this to equal $+1$ gives $\deg(m_k) = 2 - k$, in agreement with~\eqref{eq:mk-type}. Concretely: $m_1$ (the BRST differential) has degree~$+1$, $m_2$ (the binary product) has degree~$0$, and $m_3$ (the first homotopy) has degree~$-1$.
\end{proof}

\subsubsection{Rectification: from axioms to a $W(\mathsf{SC}^{\mathrm{ch,top}})$--algebra}
\label{subsec:rectification}
Conversely, suppose $A$ carries $(Q,\partial,\{m_k\})$ satisfying the axioms of Section~\ref{sec:axioms-Ainfty-chiral}.

\begin{definition}[Tameness]\label{def:tameness}
We say that an $A_\infty$ chiral algebra $(A, \{m_k\})$ is \emph{tame} if:
\begin{enumerate}[label=(\alph*)]
\item the spectral parameter dependence of $m_k$ is polynomial in $\lambda_i$ (regular part) plus a Laurent tail in $\lambda_i^{-1}$ (singular part) with finitely many negative powers;
\item the operations $m_k$ are continuous with respect to the natural topology on $A((\lambda_1)) \cdots ((\lambda_{k-1}))$;
\item for sufficiently large $k$, the operations $m_k$ vanish (truncation).
\end{enumerate}
\end{definition}

\begin{remark}[Tameness for HT theories]
\label{rem:tameness-HT}
For the $A_\infty$ chiral algebras arising from holomorphic-topological
field theories (Theorem~\ref{thm:physics-bridge}), tameness is
automatic:
\begin{enumerate}[label=\textup{(\alph*)}]
\item The spectral parameter dependence comes from meromorphic
 expansions of propagator products, which are rational functions of the
 $\lambda_i$ with finitely many poles; this gives the
 polynomial$+$Laurent structure.
\item Continuity follows from the smooth dependence of the Feynman
 integrals on the test functions and the topology of the configuration
 spaces $\FM_k(\C) \times \Conf_k^{<}(\R)$.
\item Truncation ($m_k = 0$ for $k \gg 0$) holds when the loop order of
 each Feynman diagram contributing to $m_k$ is bounded by a function
 that eventually exceeds the loop budget. Concretely, for a theory
 with $n$-valent interactions whose vertex valence strictly exceeds the
 $\mu$-slot count, the graph combinatorics force $m_k = 0$ for
 $k \gg 0$. When the $\mu$-degree equals the vertex valence (as for
 Virasoro, where each cubic vertex has exactly 2 $\mu$-slots), no such
 bound applies: diagrams of arbitrarily high loop order contribute at
 every degree, the $A_\infty$ structure is infinite
 (cf.\ Proposition~\ref{prop:vir-truncation}), and tameness
 condition~(c) must be replaced by convergence of the bar codifferential
 on the completed cofree coalgebra.
\end{enumerate}
\end{remark}

\begin{theorem}[Axioms $\Rightarrow$ operad]
\label{thm:rectification}
Assume $A$ satisfies conditions~\textup{(a)--(b)} of Definition~\ref{def:tameness}. If condition~\textup{(c)} also holds (truncation), then there exists a (canonically defined up to contractible choice) $C_\ast\!\bigl(W(\mathsf{SC}^{\mathrm{ch,top}})\bigr)$-algebra structure on~$A$ inducing the given $m_k$ by evaluation on fundamental chains as above. More generally, if condition~\textup{(c)} is replaced by convergence of the bar codifferential on the \emph{completed} reduced cofree coalgebra $\widehat{\barBch}(A)=\prod_{n\ge1}(s^{-1}\bar A)^{\otimes n}$ (as holds for the Virasoro and $W_N$ algebras), the same conclusion holds in the completed setting.
\end{theorem}

\begin{proof}
The tame \(A_\infty\) data determine an \(E_1\) chiral coalgebra over
\((\mathrm{ChirAss})^!\); the cobar functor for the two-coloured
operad \(\mathsf{SC}^{\mathrm{ch,top}}\) turns it into a strict operad
algebra whose fundamental-chain operations are the original \(m_k\).

\medskip
\textbf{Bar coalgebra from tame \(A_\infty\) data.}
The tame $A_\infty$ chiral data $(Q, \partial, \{m_k\}_{k \ge 1})$
define a codifferential on the cofree conilpotent graded coalgebra
\[
 \mathcal{C}_A \;:=\; \barBch(A)
 \;=\; \bigoplus_{n \ge 1} (s^{-1}\bar{A})^{\otimes n},
\]
where $s^{-1}$ is desuspension and $\bar{A} = A / k\cdot\mathbf{1}$
is the augmentation ideal; the coaugmented coalgebra is
$\mathrm B^{\mathrm{ch}}(A)=k\oplus\mathcal C_A$. The codifferential
$d_{\mathcal{C}} \colon \mathcal{C}_A \to \mathcal{C}_A$
is the unique coderivation whose cogenerator projection
$\pi \circ d_{\mathcal{C}} \colon \mathcal{C}_A \to s^{-1}\bar{A}$
encodes the operations $m_k$: explicitly,
\[
 \pi \circ d_{\mathcal{C}}\bigl(
 [a_1 | \cdots | a_n]\bigr)
 \;=\; \sum_{k=1}^{n}\;\sum_{i=0}^{n-k}
 (-1)^{\epsilon(i,k)}\,
 s^{-1}\,m_k(a_{i+1}, \ldots, a_{i+k}).
\]
\emph{Square-zero coderivation.} The $A_\infty$ identity
\eqref{eq:ainfty-relation-raw} is equivalent to the vanishing of
$\pi \circ d_{\mathcal{C}}^2$. Since $d_{\mathcal{C}}$ is a
coderivation on a cofree coalgebra, $d_{\mathcal{C}}^2$ is again a
coderivation; a coderivation on a cofree coalgebra is zero if and
only if its cogenerator projection is zero \cite{LV12}. Hence
$d_{\mathcal{C}}^2 = 0$.

\emph{\((\mathrm{ChirAss})^!\)-coalgebra.}
As a standalone coalgebraic object, $\mathcal{C}_A$ is an $E_1$
chiral coassociative coalgebra: the deconcatenation coproduct
encodes the $E_1$ (topological/open) colour, and the codifferential
$d_{\mathcal{C}}$ from the operations $m_k$ encodes the chiral
(holomorphic/closed) colour. Concretely, the $m_k$ take values in
iterated Laurent series
$A((\lambda_1))\cdots((\lambda_{k-1}))$, encoding the spectral
parameter dependence along $\FM_k(\C)$. The sesquilinearity axiom
\eqref{eq:sesqui-left}--\eqref{eq:sesqui-right} ensures that the
cogenerator maps intertwine the $\C$--translation action on
$\FM_k(\C)$ with $\partial$ on $A$.
The spectral substitution rule (Remark~\ref{rmk:spectral-substitution})
ensures that the cogenerator maps are compatible with the boundary
stratification of $\FM_k(\C)$: when a block $I$ collapses, the
effective spectral parameter is $\Lambda_I = \sum_{i \in I} \lambda_i$,
matching the face map of the cooperad.
Together, these equip $\mathcal{C}_A$ with the structure of a
conilpotent coalgebra over $(\mathrm{ChirAss})^!$
(the Koszul dual cooperad of the chiral associative operad),
with the spectral parameter dependence enriching it beyond the
purely topological $E_1$ bar. The full $\SCchtop$ structure lives
on the derived chiral center pair $(C^\bullet_{\mathrm{ch}}(A,A), A)$, not
on $\mathcal{C}_A$ alone: the lift from
$(\mathrm{ChirAss})^!$-coalgebra to $\mathbf B(\SCchtop)$-coalgebra
uses the promoted pair $(A,A)$ through the cobar functor on the
two-coloured operad.
Tameness (Definition~\ref{def:tameness}) is essential here:
\begin{itemize}
 \item condition (a) (polynomial$+$Laurent spectral dependence)
 ensures that each cogenerator map lands in the algebraic
 completion $A((\lambda_1))\cdots((\lambda_{k-1}))$, so the
 pairing with chains on $\FM_k(\C)$ is well-defined;
 \item condition (b) (continuity) ensures the cogenerator maps
 extend to the completed tensor products that arise from
 compositions in the bar cooperad;
 \item condition (c) (eventual truncation: $m_k = 0$ for $k \gg 0$)
 ensures that the coalgebra is concentrated in finitely many
 degrees, so the cofree coderivation is determined by finitely
 many cogenerator projections and all tree sums terminate.
\end{itemize}

\medskip
\textbf{Cobar functor and strict operad algebra.}
Apply the cobar functor $\Omega$ to $\mathcal{C}_A$:
\[
 \widetilde{A} \;:=\; \Omega(\mathcal{C}_A).
\]
By definition, $\widetilde{A}$ is the free
$\mathsf{SC}^{\mathrm{ch,top}}$--algebra on $s\,\mathcal{C}_A$
(the suspension of the $E_1$ chiral coalgebra), equipped with the
differential $d_\Omega$ induced by both the internal differential
$d_{\mathcal{C}}$ and the cooperad cocompositions. Concretely,
$\widetilde{A}$ is an algebra over
$\Omega\mathbf{B}(\mathsf{SC}^{\mathrm{ch,top}})$.

Now invoke homotopy--Koszulity --- at the associated-graded
level Theorem~\ref{thm:homotopy-Koszul}, for the coupled operad
Conjecture~\ref{conj:coupled-sc-koszul}, assumed here --- the
canonical counit
\[
 \epsilon \colon
 \Omega\mathbf{B}(\mathsf{SC}^{\mathrm{ch,top}})
 \;\xrightarrow{\;\sim\;}\;
 \mathsf{SC}^{\mathrm{ch,top}}
\]
is a quasi--isomorphism of two--colored dg operads. Therefore
restriction of structure along $\epsilon$ converts the
$\Omega\mathbf{B}(\mathsf{SC}^{\mathrm{ch,top}})$--algebra
$\widetilde{A}$ into a
$C_\ast\!\bigl(W(\mathsf{SC}^{\mathrm{ch,top}})\bigr)$--algebra,
since $\Omega\mathbf{B}(\mathsf{SC}^{\mathrm{ch,top}})$
\emph{is} the Boardman--Vogt resolution
$W(\mathsf{SC}^{\mathrm{ch,top}})$ (the cobar of the bar is
the canonical cofibrant replacement of the operad \cite{BM03}).
Denote the resulting operad algebra structure also by
$\widetilde{A}$.

\medskip
\textbf{Fundamental chains.}
The $C_\ast\!\bigl(W(\mathsf{SC}^{\mathrm{ch,top}})\bigr)$--algebra
$\widetilde{A}$ recovers the original operations $m_k$ on
fundamental chains. The cobar--bar counit
$\eta_A \colon A \to \Omega(\mathcal{C}_A) = \widetilde{A}$
embeds $A$ as the weight-$1$ component (single-vertex trees)
of the cobar complex. For the $k$--corolla $c_k$, the
structure map of $\widetilde{A}$ evaluated on the fundamental
chain $[\FM_k(\C)] \times [\mathrm{ord}]$ reduces, on
weight-$1$ elements, to the cogenerator projection of
$d_{\mathcal{C}}$:
\[
 \widetilde{m}_k(a_1, \ldots, a_k)
 \;=\; \pi \circ d_{\mathcal{C}}
 \bigl([a_1 | \cdots | a_k]\bigr)
 \;=\; m_k(a_1, \ldots, a_k).
\]
The spectral parameters are preserved because the
cogenerator maps encode exactly the Laurent coefficients of
the $\FM_k(\C)$--forms, and the bar--cobar adjunction respects
the topology on iterated Laurent series (which is where
tameness condition~(b) is used).

For composite trees $T$ with internal edges, the structure map
$\rho(T)$ of $\widetilde{A}$ is determined inductively by the
operad composition from the corolla data $\{m_k\}$. The
compatibility conditions are:
\begin{enumerate}[label=(\roman*)]
 \item The $A_\infty$ relations \eqref{eq:ainfty-relation-raw}
 ensure $d_{\mathcal{C}}^2 = 0$, which is exactly the
 statement that $\rho \circ \partial_{\FM} = d_A \circ \rho$
 on $C_\ast(\FM_k(\C))$ (boundary maps to compositions).
 \item Sesquilinearity \eqref{eq:sesqui-left}--\eqref{eq:sesqui-right}
 ensures that $\rho$ intertwines the $\C$--action on $\FM_k(\C)$
 with $\partial$ on $A$.
 \item The spectral substitution rule ensures that boundary
 faces $D_I \cong \FM_r(\C) \times \FM_{|I|}(\C)$ map
 correctly under $\rho$, with the block--sum
 $\Lambda_I = \sum_{i \in I} \lambda_i$ as effective parameter.
\end{enumerate}
These are precisely the axioms of a
$C_\ast(W(\mathsf{SC}^{\mathrm{ch,top}}))$--algebra
(cf.\ the forward direction, Proposition~\ref{prop:operad-implies-axioms}).

\medskip
\textbf{Uniqueness via Quillen equivalence.}
It remains to show that the
$C_\ast\!\bigl(W(\mathsf{SC}^{\mathrm{ch,top}})\bigr)$--algebra
structure is unique up to contractible ambiguity. By
Theorem~\ref{thm:homotopy-Koszul} and
Theorem~\ref{thm:scchtop-bar-cobar-adjunction} (for the coupled
operad, conditionally on
Conjecture~\ref{conj:coupled-sc-koszul}), the bar--cobar adjunction
\[
 \Omega \colon
 \mathrm{Coalg}_{\mathbf{B}(\mathsf{SC}^{\mathrm{ch,top}})}
 \rightleftarrows
 \mathrm{Alg}_{\mathsf{SC}^{\mathrm{ch,top}}}
 :\! \mathbf{B}
\]
is a \emph{Quillen equivalence}: the unit
$A \to \Omega\mathbf{B}(A)$ and counit
$\mathbf{B}\Omega(C) \to C$ are weak equivalences.
Suppose $\rho, \rho'$ are two
$C_\ast(W(\mathsf{SC}^{\mathrm{ch,top}}))$--algebra structures on
$A$ that agree on fundamental chains (i.e., induce the same $m_k$).
Both structures produce the same bar coalgebra $\mathcal{C}_A$
(since the bar coalgebra is determined by the cogenerator
projections, which are the $m_k$). The Quillen equivalence ensures
that the cobar $\Omega(\mathcal{C}_A)$ is the unique (up to weak
equivalence) operad algebra recovering $\mathcal{C}_A$ via the bar
construction. More precisely, the mapping space
\[
 \mathrm{Map}_{\mathrm{Alg}}^{\mathbf{B}\text{-equiv}}
 \!\bigl(\widetilde{A}, \widetilde{A}\bigr)
\]
of self-equivalences of $\widetilde{A}$ that induce the identity on
$\mathbf{B}(\widetilde{A}) \cong \mathcal{C}_A$ is contractible:
the derived unit $\eta \colon \mathrm{id} \xrightarrow{\sim}
\Omega \circ \mathbf{B}$ of the Quillen equivalence provides a
canonical contraction.
(Over a field of characteristic zero, the model category
$\mathrm{Alg}_{\mathsf{SC}^{\mathrm{ch,top}}}$ is right proper
and the mapping spaces are computed by derived hom; contractibility
of the fiber follows from the fact that $\eta_A$ is a weak
equivalence \cite{Val07, BM03}.)

\medskip
\textbf{Equivalence of categories.}
The forgetful functor $U$ (extracting $m_k$ from
$\rho([\FM_k(\C)] \times [\mathrm{ord}])$) is the construction of
Proposition~\ref{prop:operad-implies-axioms}. The bar--cobar
reconstruction above provides an inverse $R$ with
$U \circ R \cong \mathrm{id}$ on fundamental chains and
$R \circ U \simeq \mathrm{id}$ by Quillen uniqueness up to
contractible ambiguity). Hence the homotopy categories are
equivalent:
\[
 \mathrm{Ho}\bigl(\text{tame axiomatic }A_\infty\text{ chiral algebras}\bigr)
 \;\simeq\;
 \mathrm{Ho}\bigl(W(\mathsf{SC}^{\mathrm{ch,top}})\text{-algebras with
 }\partial,\,\mathbf{1}\bigr).
 \qedhere
\]
\end{proof}

\begin{remark}[Lagrangian skeleton and Stokes origin of the $A_\infty$ relations]
\label{rem:lagrangian-skeleton}
The rectification theorem (Theorem~\ref{thm:rectification}) and
its converse (Proposition~\ref{prop:operad-implies-axioms}) have a
geometric interpretation through the Lagrangian correspondence
$\mathfrak{S}_b = \mathcal{L}\times_{\mathcal{M}}\mathcal{L}$
of Remark~\ref{rem:lagrangian-geometry}.

The $A_\infty$ relations \eqref{eq:ainfty-relation-raw} are
Stokes' theorem on the boundary strata of this correspondence.
Each codimension-$1$ face of $\FM_k(\C)$ is a collision stratum
of the Lagrangian (a locus where two or more sheets of
$\mathcal{L}$ meet transversally), and the sum of residues
along these faces vanishes by Stokes, which is the Stasheff
identity $\sum m_i\circ m_j = 0$. The spectral substitution
law (Remark~\ref{rmk:spectral-substitution}) is the cluster
structure of the collision strata: when a block of points
collapses, the Lagrangian sheets merge and the effective
spectral parameter is the sum of the inner ones, reflecting
the additivity of the shifted symplectic pairing on strata.

Sesquilinearity \eqref{eq:sesqui-left}--\eqref{eq:sesqui-right}
is $\C$-equivariance of the Lagrangian embedding: the translation
action on $\FM_k(\C)$ preserves $\mathcal{L}$, and the spectral
parameters record the infinitesimal action on the normal bundle.
Tameness (Definition~\ref{def:tameness}) is the algebraicity
of the correspondence: the polynomial$+$Laurent structure
is finite-dimensionality of the fibers, and truncation is
properness of the projection.

In this language, rectification is the \emph{Lagrangian skeleton
theorem}: the self-intersection data of $\mathcal{L}$
determines the correspondence $\mathfrak{S}_b$ and hence the
$\mathsf{SC}^{\mathrm{ch,top}}$-algebra, just as the
Steinberg variety is determined by its moment-map fibers
and recovers the Hecke algebra from the convolution product
on its Borel--Moore homology.
\end{remark}

\begin{proposition}[Rectification is the Lagrangian skeleton theorem]
\label{prop:rectification-lagrangian-skeleton}
The rectification theorem \textup{(}recovering an
$\mathsf{SC}^{\mathrm{ch,top}}$-algebra from $A_\infty$ axioms\textup{)}
is the Lagrangian skeleton theorem: the self-intersection data
$\mathfrak{S}_b = \mathcal{L}\times_{\mathcal{M}}^h\mathcal{L}$
with its groupoid structure determines the Lagrangian correspondence,
and conversely. The $A_\infty$ axioms are the minimal Lagrangian
skeleton data.
\end{proposition}

\begin{proof}
By Theorem~\ref{thm:homotopy-Koszul}, the depth-zero associated
graded of the operad
$\mathsf{SC}^{\mathrm{ch,top}}$ is homotopy-Koszul (the coupled
statement is Conjecture~\ref{conj:coupled-sc-koszul}, assumed
where the coupled mixed sector is used). Combining this
with the bar-cobar adjunction (Vol~I, Theorem~A, the twisting morphism representability theorem) gives a Quillen
equivalence
\[
 \Omega\barB \colon
 \mathsf{SC}^{\mathrm{ch,top}}\text{-}\mathrm{Alg}
 \;\rightleftarrows\;
 A_\infty\text{-}\mathrm{Alg}
 \;\colon\; \barB
\]
between $\mathsf{SC}^{\mathrm{ch,top}}$-algebras and $A_\infty$
algebras satisfying the axioms of Section~\ref{sec:axioms-Ainfty-chiral}.
In the Lagrangian language: the $A_\infty$ axioms encode the
Lagrangian $\mathcal{L} \hookrightarrow \mathcal{M}$; the bar
complex $\barB$ passes to the self-intersection groupoid
$\mathfrak{S}_b = \mathcal{L}\times_{\mathcal{M}}^h\mathcal{L}$;
and the cobar functor $\Omega$ recovers the
$\mathsf{SC}^{\mathrm{ch,top}}$-algebra from the groupoid. The
Quillen equivalence states that the $A_\infty$ axioms and the
$\mathsf{SC}^{\mathrm{ch,top}}$-algebra structure determine each
other: $\barB$ passes from the SC-algebra to the $A_\infty$ data,
$\Omega$ recovers the SC-algebra from the bar coalgebra, and the
two operations are inverse on the homotopy category. In the
Lagrangian interpretation, this is the skeleton theorem: the
self-intersection data (groupoid) determines the correspondence
(SC-algebra), and conversely, because the bar-cobar adjunction
provides the independent algebraic content without circularity.
\end{proof}

\begin{remark}[Critical level exclusion discipline]
\label{rem:critical-level-exclusion-discipline}
Every $\kappa$ formula of the Sugawara form
\[
 \kappa_T(k) \;=\; \frac{\dim(\fg)\,(k+h^\vee)}{2h^\vee},
 \qquad
 r(z) \;=\; \frac{\Omega}{(k+h^\vee)\,z},
 \qquad
 T_{\mathrm{Sug}}(z) \;=\; \frac{1}{2(k+h^\vee)}\,{:}J^a J_a{:}(z),
\]
appearing throughout Volume~II (the rational-function prefactor
$(k+h^\vee)^{-1}\cdot f(k,h^\vee)$) is a rational function of $k$
with a simple pole at the \emph{critical level} $k=-h^\vee$.
At that level $T_{\mathrm{Sug}}$ degenerates: the Sugawara vector
ceases to be a conformal vector and becomes central, and the
chiral algebra $V_{-h^\vee}(\fg)$ is controlled by a
\emph{different} operative invariant, the Feigin--Frenkel centre
$\mathfrak{z}(\widehat{\fg})$ on the space of $\fg^L$-opers
(Feigin--Frenkel, Frenkel--Gaitsgory, Arakawa).

Every Sugawara-type invariant therefore carries \emph{two companion
theorems}:
\begin{enumerate}
\item[(i)] \textbf{Non-critical regime} $k\neq -h^\vee$: the
$\kappa_T$-based statement (Sugawara conformal vector,
$r(z)=\Omega/((k+h^\vee)z)$, $E_3$-topologisation,
Drinfeld strictification, spectral $R$-matrix, universal
holography).
\item[(ii)] \textbf{Critical regime} $k=-h^\vee$: the
Feigin--Frenkel-based statement. The role of the Sugawara
invariant is played by
\[
 \kappa_T^{\mathrm{FF}} \;:=\; \mathfrak{z}(\widehat{\fg})
 \;=\; \mathrm{Fun}\,\mathrm{Op}_{\fg^L}(D),
\]
an infinite-dimensional polynomial algebra on opers; the chiral
Hochschild theorem at critical level
(Theorem~\ref{thm:feigin-frenkel-chirhoch} of
Chapter~\ref{ch:hochschild}) replaces the non-critical
concentration statement by an explicit identification of
$\ChirHoch^0(V_{-h^\vee}(\fg))$ with the Feigin--Frenkel centre.
\end{enumerate}
The union of (i) and (ii) covers all levels; neither subsumes
the other. In particular, the affine $r$-matrix
$r(z)=\Omega/((k+h^\vee)z)$ diverges at $k=-h^\vee$ and must be
excluded there; the level-stripped diagnostic $r = 0$ at $k=0$ is a
fact about the rescaled $r^{\mathrm{Sug}}(z) = k\,\Omega/((k+h^\vee)z)$,
not about the Killing--Casimir $r$-matrix, and both diagnostics
coexist with exclusion of the pole at the critical level.

\emph{Scope convention adopted throughout Volume~II.} Every statement
involving the prefactor $(k+h^\vee)^{-1}$ or the Sugawara vector
$T_{\mathrm{Sug}}(z)$ carries the implicit hypothesis $k\neq -h^\vee$,
or is explicitly flagged as a rational function of~$k$ with a simple
pole at the critical level. Corollaries claimed at \emph{all} levels
produce the companion critical-level statement in terms of the
Feigin--Frenkel invariants, or restrict scope to $k\neq -h^\vee$.

\emph{Companion invariants, not exclusions.} The critical level
is not a pathological exception but the locus where a parallel
algebraic structure (the Feigin--Frenkel centre) takes over; per the
fingerprint classification, the Platonic target is the five-class
stratum $\{G,L,C,M,FF\}$ with $FF$ the canonical $\kappa=0$ companion
class.
\end{remark}

\begin{remark}[$S$-transform vs.\ Wick-rotated $R$-matrix: a chiral bridge]
\label{rem:s-transform-vs-wick-rotation-bridge}
The $S$-transform and the Wick-rotated $R$-matrix are frequently
conflated as ``the same modular datum''. They are not: they act on
different colours of $\SCchtop$ and on different algebraic strata.
\begin{enumerate}[label=(\alph*)]
\item \emph{$S$-transform (closed colour).} The $SL_2(\Z)$ generator
 $\tau \mapsto -1/\tau$ acts on (weak) modular forms and, via Zhu's
 theorem~\cite{Zhu96}, on the closed-colour chiral homology
 $H^\bullet_{\mathrm{ch}}(\Sigma_g, \A)$ at $g=1$. At a root of
 unity it descends to the Verlinde fusion endomorphism of the
 category of integrable modules
 \cite{FrenkelZhu92,HuangLepowskyKong}.
\item \emph{Wick-rotated $R$-matrix (open colour).} The spectral
 $R$-matrix $R(z) \in \End(V \otimes V)((z))$ is holomorphic on
 $\Im(z) > 0$ (time-ordered scattering) and extends by analytic
 continuation to $\Im(z) < 0$ (anti-time-ordered). The Wick
 rotation $z \mapsto iz$ exchanges the two half-planes and acts on
 the ordered $E_1$-chiral structure via the monodromy representation
 $\mathrm{Mon}(R) \colon B_n \to \End(V^{\otimes n})$.
\item \emph{Bridge at $\tau = is$.} On the Euclidean torus with
 purely imaginary modulus $\tau = is$ ($s \in \R_{>0}$), the two
 data coincide as numerical partition-function insertions:
 \[
 S \cdot Z_{\A}(\tau) \Big|_{\tau = is}
 \;=\; \mathrm{tr}_{\mathcal H} \bigl(
 \mathrm{Mon}(R)_{\mathrm{Wick}} \cdot q^{L_0 - c/24}
 \bigr) \Big|_{q = e^{-2\pi s}},
 \]
 where the right-hand side is the Wick-rotated monodromy trace
 on the ordered $E_1$ chain model and the left is the $S$-image
 of the closed-colour character. The categorical content of this
 identity is:
 \emph{$S$-transform on closed chiral homology
 $\;\equiv\;$ Wick-rotated $\mathrm{Mon}(R)$ at $\tau = is$},
 reading across the closed/open divide of $\SCchtop$ via the
 annular bar model $B^{\mathrm{ann}}(\A)$.
\end{enumerate}
\emph{Scope.} The bridge identity is PROVED for the Heisenberg
algebra (direct theta-function computation), for affine Kac--Moody
$V_k(\mathfrak g)$ at non-critical integer level (via
Frenkel--Zhu~\cite{FrenkelZhu92} and the Knizhnik--Zamolodchikov
monodromy identification), and for minimal Virasoro models
$\mathrm{Vir}_{c_{p,q}}$ (via modular-tensor-category
Verlinde/Reshetikhin--Turaev matching). It is CONJECTURAL for
non-rational chiral algebras, where both sides require separate
regularisation and the pairing of closed modular data with open
braid monodromy is not yet a theorem.
The $S$-transform (closed colour, complex structure) and the Wick
rotation of $R$ (open colour, $E_1$ ordering) are thus distinct
algebraic data that coincide numerically on the Euclidean torus
slice in the proved scope above, not universally.
\end{remark}

\subsubsection{$(H1)$--$(H5)$ as the FMP on the derived chiral algebra stack}
\label{subsec:h-axioms-from-fmp}

The five hypotheses $(H1)$--$(H5)$ used elsewhere in Volume~II
(one-loop finiteness; propagator meromorphicity and $\R$-regularity;
configuration-space convergence and AOS relations; factorisation
compatibility; conformal-vector admissibility) are not a list of
ad hoc analytic assumptions. They are the axioms of a \emph{formal
moduli problem} (FMP, in the sense of \cite[Thm.~2.2.4]{Pridham}
and \cite[Ch.~15, Thm.~1.3.12]{HA}) attached to the derived stack
$\mathbf{ChirAlg}^{\omega,\mathrm{log}}_X$ of logarithmic
$\SCchtop$-algebras with conformal vector, presented through the
factorisation ambient of Francis--Gaitsgory~\cite{FG12}.

\begin{theorem}[$(H1)$--$(H5)$ from the formal moduli problem]
\label{thm:h-axioms-from-fmp}
Let $X$ be a smooth curve over a field of characteristic~$0$ and let
\[
 \mathcal{F}_{\mathrm{log\text{-}SC}} \colon
 \mathrm{dgArt}^{\mathrm{aug}}_k \longrightarrow \mathcal{S}
\]
be the functor sending an augmented dg Artinian local algebra
$(R,\mathfrak m)$ to the space of $R$-families of
logarithmic $\SCchtop$-algebras on $X$ equipped with a
conformal vector (Definition~\ref{def:log-SC-algebra} plus a
non-critical Sugawara datum). Then:
\begin{enumerate}[label=(\roman*)]
\item \emph{FMP property.} $\mathcal{F}_{\mathrm{log\text{-}SC}}$
 is an FMP in the sense of
 \cite[\S2.1]{Pridham} / \cite[Def.~15.1.1.8]{HA}. That is,
 $\mathcal{F}_{\mathrm{log\text{-}SC}}(k) \simeq \mathrm{pt}$ and
 $\mathcal{F}_{\mathrm{log\text{-}SC}}$ preserves pullbacks along
 small surjections.
\item \emph{Tangent $L_\infty$ structure.} There is a canonical
 dg Lie (equivalently, $L_\infty$) algebra
 $\mathfrak g^{\mathrm{SC}}_X$ controlling
 $\mathcal{F}_{\mathrm{log\text{-}SC}}$ under the Lurie--Pridham
 equivalence $\mathrm{FMP}_k \simeq L_\infty\text{-}\mathrm{Alg}_k$
 \cite[\S4]{Pridham},\cite[Thm.~15.3.4.2]{HA}; explicitly
 $\mathfrak g^{\mathrm{SC}}_X$ is the convolution $L_\infty$
 algebra associated with the bar twisting cochain
 $\tau_{\mathrm{SC}} \colon B(\SCchtop) \to \SCchtop$.
\item \emph{Axiomatic coincidence.} For $(R,\mathfrak m)$ an
 augmented dg Artinian base, the enumerated hypotheses
 $(H1)$--$(H5)$ imposed on an $R$-family $(\A/R, \omega, Q, S)$
 are equivalent, one at a time, to the defining conditions
 of $\mathcal{F}_{\mathrm{log\text{-}SC}}(R)$:
 \begin{itemize}
 \item $(H1)$ \textup{one-loop finiteness} $\iff$ propagator
 factorisation on $\FM_k(\C) \times \Conf_k(\R)$ extends
 across the exceptional divisor of the Fulton--MacPherson
 compactification over~$R$, i.e.\ existence of the
 $R$-family of weight forms (tangent-obstruction vanishing
 in $H^2(\mathfrak g^{\mathrm{SC}}_X / R)$);
 \item $(H2)$ \textup{propagator regularity} $\iff$
 logarithmic-singularity condition
 $\omega_k^{\mathrm{hol}} \in \Omega^\bullet_{\log}(\FM_k(\C))$
 of Definition~\ref{def:log-SC-algebra};
 \item $(H3)$ \textup{AOS identities on configuration space}
 $\iff$ Maurer--Cartan equation for $\tau_{\mathrm{SC}}$
 on the bar cooperad $B(\SCchtop)$;
 \item $(H4)$ \textup{factorisation compatibility} $\iff$
 functoriality of $\mathcal{F}_{\mathrm{log\text{-}SC}}$
 under disjoint-union embeddings
 $\mathrm{Ran}(U) \hookrightarrow \mathrm{Ran}(X)$ in
 the chiral Koszul ambient of~\cite{FG12};
 \item $(H5)$ \textup{conformal-vector admissibility at
 non-critical level} $\iff$ choice of Sugawara lift
 $T \in \mathrm{MC}(\mathfrak g^{\mathrm{SC}}_X)$ with
 $\kappa_T \neq 0$, identifying
 $\mathcal{F}_{\mathrm{log\text{-}SC}}$ with the
 non-critical stratum of the FMP.
 \end{itemize}
\item \emph{Universality.} Every point of
 $\mathcal{F}_{\mathrm{log\text{-}SC}}(R)$, equivalently every
 Maurer--Cartan element of
 $\mathfrak g^{\mathrm{SC}}_X \otimes \mathfrak m$, is a
 logarithmic $\SCchtop$-algebra over $R$ satisfying
 $(H1)$--$(H5)$; conversely every such algebra arises this way.
\end{enumerate}
\end{theorem}

\begin{proof}[Sketch]
\emph{(i)} The pullback property for $\mathcal{F}_{\mathrm{log\text{-}SC}}$
reduces to the pullback property for the functor classifying
$C_\ast(W(\SCchtop))$-algebras, which holds because
$C_\ast(W(\SCchtop))$ is a cofibrant dg operad (Koszul resolution
of $\SCchtop$ via the coloured Koszul duality of
\cite{FG12,HA}) and algebras over cofibrant operads satisfy the
pullback axiom for FMPs.

\emph{(ii)} Apply the Lurie--Pridham correspondence
$\mathrm{FMP}_k \simeq L_\infty\text{-}\mathrm{Alg}_k$
\cite[Thm.~2.2.4]{Pridham},\cite[Thm.~15.3.4.2]{HA}. The
tangent complex of $\mathcal{F}_{\mathrm{log\text{-}SC}}$ at the
augmentation $k \to k$ is, by operadic deformation theory, the
convolution complex
$\mathfrak g^{\mathrm{SC}}_X =
\mathrm{Conv}(B(\SCchtop), \mathcal{E}nd_\A^{\mathrm{ch}})$.
Its $L_\infty$ structure is the standard convolution structure
of \cite{LV12}. Chiralisation (the extension of the shuffle
operad ambient to the BD chiral pseudo-tensor on $\mathrm{Ran}(X)$)
is precisely the Francis--Gaitsgory \cite{FG12} replacement of
$\otimes$ by $\otimes^{\mathrm{ch}}$; the resulting
$L_\infty$-algebra in $\mathrm{IndCoh}(\mathrm{Ran}(X))$
classifies the $R$-family FMP.

\emph{(iii)} The axiomatic coincidence is read off by comparing
the hypotheses of Theorem~\ref{thm:physics-bridge} (which
produces logarithmic $\SCchtop$-algebras from BV data) with the
Maurer--Cartan equation on
$\mathfrak g^{\mathrm{SC}}_X \otimes \mathfrak m$. Each
hypothesis corresponds to one component of the defining datum
(weight form existence, logarithmic-pole condition, MC equation,
Ran-descent, Sugawara lift).

\emph{(iv)} Universality is the content of the Lurie--Pridham
correspondence: MC elements of $\mathfrak g^{\mathrm{SC}}_X
\otimes \mathfrak m$ and $R$-points of the FMP are in
equivalence.
\end{proof}

\begin{corollary}[$(H1)$--$(H5)$ are not ad hoc]
\label{cor:h-axioms-universal}
The hypotheses $(H1)$--$(H5)$ are equivalent to membership in the
universal FMP $\mathcal{F}_{\mathrm{log\text{-}SC}}$ classifying
deformations of logarithmic $\SCchtop$-algebras with conformal
vector. They are not independent analytic assumptions imposed by
the construction; they are the component-wise unpacking of a
single derived-geometric structure (the tangent $L_\infty$ algebra
$\mathfrak g^{\mathrm{SC}}_X$).
\end{corollary}

\begin{remark}[Scope]
\label{rem:fmp-scope}
Theorem~\ref{thm:h-axioms-from-fmp}(i)--(iv) is proved
unconditionally in the classical limit of the chiral ambient,
i.e.\ over the formal disk $D = \Spec k[[z]]$, where
$\mathfrak g^{\mathrm{SC}}_D$ is a genuine $L_\infty$-algebra over
$k$ and the Lurie--Pridham correspondence applies directly. For
the global curve $X$ and, more generally, for families over the
moduli stack $\overline{\mathcal M}_{g,n}$, the tangent datum is
the $L_\infty$-algebra
$\mathfrak g^{\mathrm{SC},\mathrm{mod}}_X$ whose MC elements
encode modular compatibility; the reduction of $(H1)$--$(H5)$ to
FMP structure over $\overline{\mathcal M}_{g,n}$ coincides with
the modular bootstrap (curved Dunn additivity at
$g\ge 1$, $\kappa_T$-weight form compatibility across stable
graphs, cf.~Thm.~\ref{thm:curved-dunn-H2-vanishing-all-genera}).
\end{remark}

\subsubsection{Compatibility with the PVA construction}
\label{subsec:compat-PVA}

The axiomatic and operadic incarnations of an $\Ainf$ chiral
algebra having been shown equivalent, it remains to verify that
the PVA structure produced on cohomology
(Theorem~\ref{thm:cohomology_PVA}) agrees with the one read
directly off the associated graded of the operad. This is a
consistency check, not a new construction.

The proof of the PVA structure on cohomology in
Section~\ref{sec:axioms-Ainfty-chiral} can be interpreted
operadically as the statement that the associated graded
(holomorphic filtration) of $W(\mathsf{SC}^{\mathrm{ch,top}})$ is
the free product of BD chiral and $E_1$, so the cohomological
operations descend from the closed color; the bracket is the
classical limit (residue) of the operadic composites, and
associativity is the regular part of the $m_2$ identity.

\begin{remark}[Good Z-grading scope for DS reduction]\label{rem:ds-good-grading-scope}
Quantum Drinfeld--Sokolov reduction $\mathcal{W}_k(\mathfrak{g}, f)$ requires a \emph{good} $\mathbb{Z}$-grading on $\mathfrak{g}$ with $f \in \mathfrak{g}_{-2}$: $\operatorname{ad} f \colon \mathfrak{g}_0 \to \mathfrak{g}_{-2}$ surjective, Kazhdan pairing on $\mathfrak{g}_{1/2}$ non-degenerate, $f$ in a unique nilpotent orbit meeting $\mathfrak{g}_{-2}$. The principal nilpotent always admits a good grading. Classical types $A_n, B_n, C_n, D_n$ admit good gradings on their nilpotent orbits (Kac--Roan--Wakimoto); in exceptional types $E_{6,7,8}, F_4, G_2$ the admissible orbits are the classified good-grading orbits. This is the scope for forming the DS-BRST complex. It is not a DS--Hochschild transport theorem, not a chain-level $\SCchtop$ lift, and not a Koszul-preservation statement. In this volume the DS--Hochschild comparison is proved for principal and hook-type data, with named cover-descent fibres such as Bershadsky--Polyakov checked separately; subregular, minimal beyond named fibres, exceptional good-graded, and exotic nilpotents require case-by-case BRST admissibility, invariant transferred braces, and Hochschild comparison.
\end{remark}

\begin{remark}[Five $E_1$-chiral notions and the present axioms]
\label{rem:axioms-five-e1-notions}%
\index{E_1-chiral algebra@$\Eone$-chiral algebra!five notions}%
The single phrase ``$\Eone$-chiral algebra'' names at least five
distinct objects in the literature, none of them interchangeable
with another at chain level; the present axiomatic notion must be
anchored canonically among them.
(A) Strict $\mathrm{ChirAss}$-algebra in $\mathrm{D}\text{-}\mathsf{mod}(X)$
(Beilinson--Drinfeld pseudo-tensor): the bare strict associative
chiral algebra.
(B) $A_\infty$-algebra in $\End^{\mathrm{ch}}_\cA$: the present
Definition~\ref{def:ainfty_chiral} together with the
$\Eone$-chiral interpretation of
Definition~\ref{def:e1-chiral-algebra}; equivalently, an
$A_\infty$-algebra structure on the chiral endomorphism operad
of Remark~\ref{rem:operadic-formulation}.
(C) Etingof--Kazhdan quantum vertex algebra: the Drinfeld-associator-fixed
quantization of an $\Einf$-chiral seed via the EK braided
operad isomorphism, well-defined on the chirally Koszul locus.
(D) $A_\infty$-algebra in the $\Eone$-chiral $\infty$-operad:
homotopy-coherent $\Eone$-algebra in $\Fact^{\mathrm{ord}}(X)$,
native to the Gaitsgory--Rozenblyum $(\infty,2)$-categorical
ambient.
(E) Factorization algebra on $\Ran^{\mathrm{ord}}(X)$: the
Francis--Gaitsgory geometric incarnation, with $n$-ary operations
on $\Conf_n^{\mathrm{ord}}(X)$ not required to descend to
$\Sigma_n$-coinvariants.
The present axioms produce (B);
(B) $\leftrightarrow$ (C) on the Koszul locus via the EK quantization
isomorphism modulo the choice of Drinfeld associator
$\Phi \in \GRT_1(\Q)$; (B) $\leftrightarrow$ (D) on the same
locus via the rectification of $A_\infty$-algebras to algebras
over a cofibrant operad (Theorem~\ref{thm:rectification});
(D) $\leftrightarrow$ (E) by Francis--Gaitsgory factorization
recognition; (A) is \emph{strictly weaker} than (B)--(E) (no
homotopy data, no associator). The downstream chapters use (B)
on the formal disk and (D)/(E) on the global curve; no statement
requires (A) alone, and every statement mentioning (C) carries
an explicit Drinfeld-associator choice. The bridges are theorems,
not definitions.
\end{remark}

\begin{remark}[The axioms as input to the universal-holography theorem]
\label{rem:axioms-holography-anchor}%
\index{universal holography!axioms anchor}%
The axiomatic boundary algebra of
Definition~\ref{def:ainfty_chiral} is the family-by-family input
to the universal-holography master theorem
(Theorem~\ref{thm:universal-holography-master}).
Each rung specialises this axiomatic input:
the Heisenberg corollary (Corollary~\ref{cor:rung-heisenberg})
takes $\cA = \cH_k$ with $\kappaChHodge(\cH_k) = k$, abelian Sugawara
$T = \tfrac{1}{2k}{:}JJ{:}$, $c = 1$;
the affine Kac--Moody corollary
(Corollary~\ref{cor:rung-affine-km}) takes $\cA = V_k(\fg)$ with
$\kappaChHodge(V_k(\fg)) = \dim(\fg)(k+h^\vee)/(2h^\vee)$ and Sugawara
stress at non-critical level;
the Virasoro corollary (Corollary~\ref{cor:rung-virasoro})
takes $\cA = \mathrm{Vir}_c$ with
$\kappaChHodge(\mathrm{Vir}_c) = c/2$ and BRST-primitive $G^{(2)}$
from the Costello--Gaiotto holomorphic Chern--Simons construction;
the $W_N$ corollary (Corollary~\ref{cor:rung-w-N}) takes
$\cA = \cW_N(\fg)$ with $\kappaChHodge(\cW_N) = c(H_N - 1)$, where
$H_N = \sum_{j=1}^N 1/j$ (boundary check at $N=2$:
$\kappaChHodge(\cW_2) = c/2 = \kappaChHodge(\mathrm{Vir})$); the
$W_\infty$ corollary (Corollary~\ref{cor:rung-w-infty}) is the
iterated-Sugawara $E_\infty$-topological limit
(Theorem~\ref{thm:e-infinity-topologization-ladder}). The
$\Sigma_n$-coinvariant projection of the ordered axioms is the
symmetric-bar shadow used in the modular-characteristic
computation; the ordered axioms themselves carry the full
$R$-matrix data via the deconcatenation coproduct.

The Volume~I Koszul-reflection identity, on the genus-$0$
affine/tangent KZ window,
$\bar\partial = \KZ^*(\nabla_{\mathrm{Arnold}})$ together with
$\kappaChHodge(\cA) = -c_{\mathrm{ghost}}(\mathrm{BRST})$ is, on the
present axiomatic side, the comparison between the bar
differential $d_{\barB}$ on $T^c(s^{-1}\bar\cA)$ assembled from
the $m_k$ above and the holomorphic Knizhnik--Zamolodchikov
connection on $\Conf_n(\C)$ pulled back via the collision-residue
Arnold relations on $\FM_n(\C)$. The universal conductor formula
$K(\cA) = \kappaChHodge(\cA) + \kappaChHodge(\cA^!)$ and the universal
trace identity bridge are the characteristic-class shadows of
the present axiomatic data under the open/closed pentagon of
equivalences for $\SCchtop$.
\end{remark}

\begin{remark}[Curved-$A_\infty$ axiomatic support for the K3 BKM
 Borcherds correction]
\label{rem:axioms-k3-bkm-curved-ainfty}%
\index{K3 chiral bialgebra!curved-$A_\infty$ support}%
\index{Borcherds-imaginary-root relation!axiomatic handle}%
\index{Hochschild concentration!CY-$3$ extension for K3}%
The K3 chiral bialgebra target $\cA = \mathbf{H}_{\Delta_5}$ of
Vol~III, Ch.\,\ref{V3-ch:k3-chiral-bialgebra-platonic}, is compared
with the axiomatic framework above only through the curved
$A_\infty$ extension (Definition~\ref{def:curved-ainfty-chiral}) and
after the finite Hall/CoHA source, pairing, PBW/no-extra-relations,
radical, parity, completion, inverse-limit, and EK/Heegner gates. In
that conditional comparison the chain-level generators realise
$V\subset\cA$ of homogeneous
OPE-weight~$1$, the quadratic Miki commutators realise $m_2$
purely in $V\otimes V$, and the Borcherds-imaginary-root relations
$R_{\mathrm{Borcherds}}$
\textup{(}Vol~I, Ch.\,\ref{V1-chap:climax-platonic},
Remark~\textup{\ref{V1-rem:k3-koszul-structure-precise}}\textup{)}
realise as higher-arity operations $m_k$ with $k\geq 3$ and as
a non-vanishing curvature element
\[
 m_0 \;=\; \omega_{\Delta_5} \;\in\; \Gamma\!\bigl(
 \overline{\mathcal{A}_2},\, \mathcal{O}_{\overline{\mathcal{A}_2}}\bigr),
 \qquad
 \deg \omega_{\Delta_5} = 2,
\]
whose EK/Heegner comparison target is the Gritsenko--Nikulin
denominator cocycle read automorphically. The $n=1$ curved relation of
Definition~\ref{def:curved-ainfty-chiral}, $m_1^2 = [m_0, -]$,
encodes the conditional test that the bar differential on the K3 input
fails to square to zero by the Humbert-divisor obstruction: on
the Koszul locus $\overline{\mathcal{A}_2}\setminus H_1$ the
curvature $m_0$ vanishes to second order in local coordinates and
the ordinary $A_\infty$ relations hold; across $H_1$ the order-$8$
monodromy of Vol~III,
Theorem~\textup{\ref{thm:humbert-order-K-kappa}}, is the
cohomological shadow of non-vanishing $m_0$.

\smallskip
\noindent\emph{The K3 categorical and chiral supports.}
Volume~I Theorem~H begins with a family datum
$H_H(\mathbf H_{\Delta_5};S_{\Delta_5})$ and gives
\[
 \ChirHoch^n(\mathbf H_{\Delta_5})
 \cong H^n(K_{\mathbf H_{\Delta_5},S_{\Delta_5}}),
 \qquad
 \operatorname{Supp}\ChirHoch^\bullet(\mathbf H_{\Delta_5})
 \subseteq S_{\Delta_5}.
\]
The HKR theorem computes the categorical Hochschild object of
$D^b(\mathrm{Coh}(K3\times E))$ in its CY$_3$ ambient.  Its passage
to the chiral curve chart requires the comparison
\[
 q_{\Phi,C}\colon
 \HH^\bullet(D^b\mathrm{Coh}(K3\times E))longrightarrow
 \ChirHoch^\bullet(\mathbf H_{\Delta_5}).
\]
The K3 comparison problem is to construct $q_{\Phi,C}$ and the
cochain homotopies that send every categorical degree outside
$S_{\Delta_5}$ to an exact chart cochain.  The degree-three chart
group is then determined by the resulting complex.  The
Schiffmann--Vasserot degree-three CoHA generator supplies a candidate
source class whose image remains to be computed.

\smallskip
\noindent\emph{Bar-cobar inversion in the curved lane.}
On the Koszul locus $\overline{\mathcal{A}_2}\setminus H_1$, the
bar-cobar counit of Theorem~A, strengthened in the completed
Positselski Theorem~B \textup{(}Vol~I,
Theorem~\textup{\ref{V1-thm:theorem-b-platonic-pointer}}\textup{)}
is the comparison target for $\mathbf{H}_{\Delta_5}$ in the
curved-$A_\infty$ category with the curved bar functor of
Positselski\,2011, Memoirs AMS\,212, Ch.~$7$; the
weight-completed coderived-contraderived equivalence must survive the
$m_0\ne 0$ deformation because
$\omega_{\Delta_5}$ is a scalar on the base and does not couple to
the bar-length filtration. The chain-level chain-homotopy witness
is the $h_{\mathrm{bar}}^{\Delta_5}$ of
Remark~\textup{\ref{V1-rem:k3-koszul-structure-precise}}, valued in the
Positselski curved CDG bar resolution.
\end{remark}

\section{K3 base-point axioms}
\label{sec:axioms-supplement}

\begin{remark}[Vol~II axioms of $\mathsf{SC}^{\mathrm{ch,top}}$ at the K3 base point]
\label{rem:axioms-1}
The Vol~II axioms of $\mathsf{SC}^{\mathrm{ch,top}}$ ($3$D HT QFT)
instantiate at the K3 base point with the boundary chart
$\mathbf H_{\Delta_5}$ at the Schur sector of class-$\mathcal S$ $A_1$
on $\Sigma_{0,24}$, with Shapere--Tachikawa central charges
$c_{4d} = 107/6$ and $c_{2d} = -214$,
the chiral base $\mathrm{Ran}(E^{\mathrm{nod,sm}}_{24})$, and the
parameter base $\bar{\mathcal A}_2$. The averaging morphism is
\[
\mathrm{av}
=\mathrm{Torelli}_{K3}\circ\mathrm{KugaSatake}\circ j_{\mathrm{Kod}}.
\]
The holomorphic-twist axiom gives the chiral boundary algebra; the
Borcherds--Gritsenko lift gives the automorphic scalar
$\Delta_5$. These are linked by $\mathrm{av}$, not identified as
the same object \cite{CostelloHolomorphic,Gaiotto2009}.
Compare Vol.~III Chapter~\ref{V3-ch:k3-chiral-bialgebra-platonic}.
\end{remark}

\begin{remark}[K3 base-point refinement]
\label{rem:axioms-2}
At the K3 base point, the Vol~II axioms refine as follows. The
factorisation axiom is bi-based: Beilinson--Drinfeld factorisation
lives on $\mathrm{Ran}(E^{\mathrm{nod,sm}}_{24})$, while the
automorphic line $\mathcal L^{\Delta_5}$ lives over
$\bar{\mathcal A}_2$. The Borcherds scalar is
$\Delta_5$ on the Maass spin cover, with
$\Phi_{10}^{\mathrm{un}}=\Delta_5^2$ on the untwisted Siegel
cover. The Hall--Drinfeld-double axiom specifies the conditional
target
\[
\mathcal D_\hbar^{\mathrm{red},\mathrm{comp}}(
  \mathcal Y^{\mathrm{Hall}}_\hbar,
  \widetilde\Phi^{\mathrm{Sieg\text{-}Bor}}_{\mathrm{Sp}_4},
  R_{\mathrm{Sieg,dyn}})
\;\dashrightarrow\;
\mathbf H_{\Delta_5},
\]
after the finite Hall/CoHA source, pairing, PBW/no-extra-relations,
radical, parity, completion, inverse-limit, and EK/Heegner comparison
gates. The $R$-matrix axiom supplies the comparator
$R_{\mathrm{Sieg,dyn}}$ via the Pasol--Zagier
Kronecker--Eisenstein series. The conductor identity
\[
\hbar_{\mathrm{QG}}^2 K^{\kappaChHodge}=-1,\qquad
K^{\kappaChHodge}=2c_+(\mathrm{Mukai}(K3))=8,
\]
fixes the quantum-group specialisation
$\hbar_{\mathrm{QG}}^2=-1/8$ \cite{PasolZagier2013,EtingofKazhdan2007PartV}.
Compare Chapter~\ref{V3-ch:k3-chiral-bialgebra-platonic}.
\end{remark}
